\documentclass[11pt]{article}

\usepackage[T1]{fontenc}
\usepackage[utf8]{inputenc}
\usepackage{lmodern}
\usepackage{microtype}
\usepackage[margin=0.85in]{geometry}
\usepackage{amsmath,amssymb,amsthm,mathtools}
\usepackage{booktabs,array,tabularx}
\usepackage{tikz-cd}
\usepackage{enumitem}
\usepackage{needspace}
\usepackage[hidelinks]{hyperref}

\newcommand{\llbracket}{\mathopen{[\mkern-3mu[}}
\newcommand{\rrbracket}{\mathclose{]\mkern-3mu]}}
\newcommand{\atsign}{\mathord{@}}
\newcommand{\finord}[1]{\underline{#1}}
\newcommand{\DP}{Do\v{s}en--Petri\'{c}}
\newcommand{\tightlist}{\setlength{\itemsep}{0pt}%
  \setlength{\parskip}{0pt}}

% Arithmetic notation used in Appendix A.
\newcommand{\N}{\mathbb N}
\newcommand{\LA}{L_{\mathrm A}}
\newcommand{\pair}[2]{\pi\bigl(#1,#2\bigr)}
\newcommand{\ctag}[2]{c\bigl(#1,#2\bigr)}
\newcommand{\PA}{\mathrm{PA}^{*}}
\newcommand{\num}{\operatorname{num}}
\newcommand{\Num}{\operatorname{Num}}
\newcommand{\scod}{\operatorname{sc}}
\newcommand{\Sc}{\operatorname{Sc}}
\newcommand{\sbop}{\operatorname{sb}}
\newcommand{\Sb}{\operatorname{Sb}}
\newcommand{\Prf}{\operatorname{Prf}}
\newcommand{\Prov}{\operatorname{Pr}}
\newcommand{\Con}{\operatorname{Con}}
\newcommand{\FV}{\operatorname{FV}}
\newcommand{\Q}{\mathsf Q}
\newcommand{\h}{\mathsf h}
\newcommand{\code}[1]{\#\bigl(#1\bigr)}
\newcommand{\quo}[1]{\ulcorner #1\urcorner}
\newcommand{\bn}[1]{\overline{#1}}
\newcommand{\flatof}[1]{\bigl(#1\bigr)^{\flat}}
\newcommand{\defeq}{\mathrel{:=}}
\newtheorem{proposition}{Proposition}[section]
\theoremstyle{definition}
\newtheorem{example}[proposition]{Example}

\setlist[itemize]{leftmargin=*,topsep=0.5\baselineskip}
\setlength{\emergencystretch}{3em}
\allowdisplaybreaks

\hypersetup{
  pdftitle={The Unbearable Notation of Logic},
  pdfsubject={Categorical structure and at-sign calculation for first-order terms, formulas, substitution, evaluation, and language expansion},
  pdfkeywords={first-order logic, term monad, Kleisli category, Eilenberg--Moore algebra, Dosen--Petric substitution, occurrence maps, Lawvere theory, classifying topos, at-sign notation, evaluation pairing, arithmetization, Feferman, numeral substitution}
}

\title{The Unbearable Notation of Logic}
%\subtitle{Categorical Structure and the \texorpdfstring{$\atsign$}{@}-Calculus}
\author{Florian Lengyel}
%\email{flengyel@gradcenter.cuny.edu}
\date{September 26, 2026}

\begin{document}
\maketitle

\begin{abstract}
Standard presentations of first-order syntax and semantics routinely use
similar brackets, superscripts, and substitution expressions for several
differently typed operations. Fix a first-order language \(\mathcal L\) with
functional signature \(\Sigma\), sets \(X,Y\) of variables, and an
\(\mathcal L\)-structure with domain \(A\). Four relevant types are:
\[
X\to T_\Sigma Y,\qquad X\to A,\qquad
T_\Sigma X\to A,\qquad
\operatorname{Form}_{\mathcal L}(X)\times A^X\to\mathbf 2.
\]
Here \(A^X\) is the set of functions \(X\to A\), and
\(\mathbf2=\{\bot,\top\}\) is the set of truth values. The four displayed maps
represent syntactic substitution, assignment, term evaluation under an
assignment, and satisfaction. Readers must ordinarily infer from convention
which operation an expression denotes.

This note presents two complementary descriptions of the operations. The categorical
description organizes terms into a monad, substitutions into Kleisli arrows,
and functional structures into Eilenberg--Moore algebras. The $\atsign$-calculus
uses prefix $\atsign$ for term evaluation and postfix $\atsign$ for satisfaction;
it is the working notation for calculations. The categorical notation records
compositional structure, while the $\atsign$-notation performs the corresponding
term-evaluation and satisfaction calculations. Adjoining names for the elements of a structure changes the
term functor from $T_\Sigma X$ to $T_\Sigma(A+X)$ and turns assignments into
substitutions by ground terms. A finite-occurrence construction of \DP\ records
how generalized composition routes repeated formula-pattern occurrences.
Functorial semantics then extends the categorical
description to internal models and, for geometric theories, to models as
generalized points of a classifying topos. An appendix separates
expressions, codes, and numeral-valued terms in Feferman's
arithmetization conventions.
\end{abstract}

First-order syntax and semantics involve terms, syntactic substitutions,
assignments, term evaluation, formulas, and satisfaction. Textbook notation
often leaves the types distinguishing those six ingredients to convention.
The term-monad description below explains how the ingredients compose; the
$\atsign$-calculus makes their substitution laws concise enough for routine
calculation. The categorical diagrams state the structural laws, and the
$\atsign$-identities state their calculational forms.

The note uses the familiar inductive definitions of terms, formulas,
evaluation, and satisfaction. Categorical notation records their compositional
laws, and the $\atsign$-calculus supplies their working calculation notation.
A later section adds the \DP\ occurrence maps, which record how
formula-pattern substitution routes the individual occurrences of the outer
pattern's formal arguments.

\section{Categorical vocabulary}\label{categorical-vocabulary}

A \textbf{category} consists of objects and arrows between them. For every
object \(X\) there is an identity arrow \(1_X:X\to X\), and composable arrows
\(f:X\to Y\) and \(g:Y\to Z\) have a composite
\(g\circ f:X\to Z\). Composition is associative, and the identity laws are
\(f\circ1_X=f=1_Y\circ f\). The category \(\mathbf{Set}\) has sets as objects
and functions as arrows.

A \textbf{functor} \(F:\mathcal C\to\mathcal D\) sends objects to
objects and arrows to arrows while preserving identities and
composition. An \textbf{endofunctor} is a functor from a category to
itself.

\Needspace{8\baselineskip}
Given functors \(F,G:\mathcal C\to\mathcal D\), a \textbf{natural
transformation} \(\theta:F\Rightarrow G\) assigns an arrow
\(\theta_X:FX\to GX\) to every object \(X\) so that, for every
\(u:X\to Y\), the following square commutes:

\[
\begin{tikzcd}[row sep=large,column sep=large]
F X \arrow[r,"\theta_X"] \arrow[d,"F(u)"'] &
G X \arrow[d,"G(u)"] \\
F Y \arrow[r,"\theta_Y"'] &
G Y
\end{tikzcd}
\]

Equivalently, \(G(u)\circ\theta_X=\theta_Y\circ F(u)\): applying
\(\theta\) is compatible with changing \(X\) by \(u\). The unit of the term
monad, defined in Section~\ref{the-term-monad}, is a natural transformation.
Its naturality says that mapping variable leaves along \(u:X\to Y\) commutes
with regarding those variables as terms.

A \textbf{monad} on a category \(\mathcal C\) is an endofunctor
\(T:\mathcal C\to\mathcal C\) together with natural transformations
\(\eta:1_{\mathcal C}\Rightarrow T\) and \(\mu:T^2\Rightarrow T\), called
the unit and multiplication. Here \(1_{\mathcal C}\) is the identity functor,
which fixes every object and arrow, while \(T^2=T\circ T\) and
\(T^3=T\circ T\circ T\) are iterated applications of \(T\). The monad laws
are typed by the following commutative diagrams. For every object \(X\), the
two unit triangles are

\[
\begin{tikzcd}[row sep=large,column sep=large]
T X \arrow[r,"\eta_{T X}"] \arrow[dr,"1_{T X}"'] &
T^2X \arrow[d,"\mu_X"] \\
& T X
\end{tikzcd}
\qquad
\begin{tikzcd}[row sep=large,column sep=large]
T X \arrow[r,"T(\eta_X)"] \arrow[dr,"1_{T X}"'] &
T^2X \arrow[d,"\mu_X"] \\
& T X,
\end{tikzcd}
\]

and associativity is the square

\[
\begin{tikzcd}[row sep=large,column sep=large]
T^3X \arrow[r,"T(\mu_X)"] \arrow[d,"\mu_{T X}"'] &
T^2X \arrow[d,"\mu_X"] \\
T^2X \arrow[r,"\mu_X"'] & T X.
\end{tikzcd}
\]

For the term monad below, \(T\) forms terms, \(\eta\) regards a variable as a
term, and \(\mu\) removes one formal term layer. Section~\ref{the-term-monad}
defines this operation recursively and computes it on an explicit term of
terms.

The \textbf{Kleisli category} \(\mathbf{Kl}(T)\) of a monad \(T\) has
the same objects as \(\mathcal C\). An arrow from \(X\) to \(Y\) in
\(\mathbf{Kl}(T)\) is an ordinary arrow \(X\to TY\) in \(\mathcal C\).
The identity on \(X\) is \(\eta_X:X\to TX\). If
\(\sigma:X\to TY\) and \(\tau:Y\to TZ\), their Kleisli composite is
\(\tau\star\sigma=\mu_Z\circ T(\tau)\circ\sigma:X\to TZ\). For the term
monad, Kleisli arrows are precisely simultaneous syntactic substitutions.

An \textbf{Eilenberg--Moore algebra} for a monad \(T\) is an object
\(A\) equipped with an arrow \(a:TA\to A\) for which the following unit
triangle and multiplication square commute:

\[
\begin{tikzcd}[row sep=large,column sep=large]
A \arrow[r,"\eta_A"] \arrow[dr,"1_A"'] & T A \arrow[d,"a"] \\
& A
\end{tikzcd}
\qquad
\begin{tikzcd}[row sep=large,column sep=large]
T^2A \arrow[r,"T(a)"] \arrow[d,"\mu_A"'] &
T A \arrow[d,"a"] \\
T A \arrow[r,"a"'] & A.
\end{tikzcd}
\]

It specifies a coherent way to interpret or evaluate the formal expressions
collected by \(T\). A homomorphism \(h:(A,a)\to(B,b)\) of such algebras is
an arrow \(h:A\to B\) for which the following square commutes:

\[
\begin{tikzcd}[row sep=large,column sep=large]
T A \arrow[r,"T(h)"] \arrow[d,"a"'] & T B \arrow[d,"b"] \\
A \arrow[r,"h"'] & B
\end{tikzcd}
\]

The algebras and their homomorphisms form the
\textbf{Eilenberg--Moore category} \(\mathcal C^T\). For the term monad, its
objects are the ordinary algebras interpreting the function symbols of a
signature.

For monads \(T,S\) on the same category, a \textbf{monad morphism}
\(j:T\Rightarrow S\) is a natural transformation for which the following
unit triangle and multiplication rectangle commute for every \(X\):

\[
\begin{tikzcd}[row sep=large,column sep=large]
X \arrow[r,"\eta_X^T"] \arrow[dr,"\eta_X^S"'] &
T X \arrow[d,"j_X"] \\
& S X
\end{tikzcd}
\]

\[
\begin{tikzcd}[row sep=large,column sep=large]
T^2X \arrow[r,"j_{T X}"] \arrow[d,"\mu_X^T"'] &
S T X \arrow[r,"S(j_X)"] & S^2X \arrow[d,"\mu_X^S"] \\
T X \arrow[rr,"j_X"'] && S X.
\end{tikzcd}
\]

Thus \(j\) translates \(T\)-expressions into \(S\)-expressions while
preserving variables and substitution.

\section{Languages and their functional
signatures}\label{languages-and-their-functional-signatures}

Let \(\mathcal L\) be a single-sorted first-order language. Separate its
nonlogical symbols into:

\begin{itemize}
\tightlist
\item
  a finitary functional signature \(\Sigma\), including constant symbols
  as nullary function symbols; and
\item
  a family \(\mathcal R\) of relation symbols with specified arities.
\end{itemize}

The term construction depends only on \(\Sigma\). Relation symbols enter
later, when formulas and satisfaction are defined.

For each set \(X\) of variables, let \(T_\Sigma X\) be the set of
\(\Sigma\)-terms whose variables belong to \(X\). Thus \(T_\Sigma X\) is
generated by variables \(x\in X\) and expressions
\(f(t_1,\ldots,t_n)\), where \(f\in\Sigma\) has arity \(n\) and each
\(t_i\) belongs to \(T_\Sigma X\).

A map \(u:X\to Y\) sends each variable leaf \(x\) to the variable leaf
\(u(x)\) and induces \(T_\Sigma u:T_\Sigma X\to T_\Sigma Y\), recursively
defined by

\[
T_\Sigma u(\eta_X(x))=\eta_Y(u(x)),\qquad
T_\Sigma u\bigl(f(t_1,\ldots,t_n)\bigr)
=f\bigl(T_\Sigma u(t_1),\ldots,T_\Sigma u(t_n)\bigr).
\]

For every \(v:Y\to Z\), structural induction yields
\(T_\Sigma(1_X)=1_{T_\Sigma X}\) and
\(T_\Sigma(v\circ u)=T_\Sigma v\circ T_\Sigma u\). Hence \(T_\Sigma\) is an
endofunctor on \(\mathbf{Set}\).

\section{The term monad}\label{the-term-monad}

The endofunctor \(T_\Sigma\) has a canonical monad structure. Its unit
\(\eta_X:X\to T_\Sigma X\) sends a variable to the corresponding variable
term. To define its multiplication, first observe that
\(T_\Sigma^2X=T_\Sigma(T_\Sigma X)\). An element of this set is an outer
\(\Sigma\)-term whose variable labels are themselves elements of
\(T_\Sigma X\). If \(p\in T_\Sigma X\), then
\(\eta_{T_\Sigma X}(p)\) is the outer variable labelled by the entire inner
term \(p\).

The multiplication \(\mu_X:T_\Sigma(T_\Sigma X)\to T_\Sigma X\) is defined
by structural recursion:

\[
\begin{aligned}
\mu_X\bigl(\eta_{T_\Sigma X}(p)\bigr)&=p,
&&p\in T_\Sigma X,\\
\mu_X\bigl(f(U_1,\ldots,U_n)\bigr)
&=f\bigl(\mu_X(U_1),\ldots,\mu_X(U_n)\bigr),
&&f\in\Sigma\text{ of arity }n.
\end{aligned}
\]

Every term has exactly one outermost constructor: it is either a variable term
or an expression \(f(U_1,\ldots,U_n)\). The two clauses therefore assign
exactly one value to every element of \(T_\Sigma(T_\Sigma X)\), and structural
recursion yields a function \(\mu_X\) with the displayed type. The operation is
often called \emph{flattening}: it replaces every outer variable by the inner
term labelling that variable and preserves every operation symbol.

\paragraph{Example: flattening a term of terms.}
Use a signature containing the nullary symbol \(1\), the unary successor
symbol \(S\), and the binary symbols \(+\) and \(\cdot\), and let
\(X=\{x,y\}\). In
\(T_\Sigma X\), put \(p=x+1\), \(q=S(y)\), and \(r=x\cdot y\), and abbreviate

\[
\widehat p=\eta_{T_\Sigma X}(p),\qquad
\widehat q=\eta_{T_\Sigma X}(q),\qquad
\widehat r=\eta_{T_\Sigma X}(r).
\]

The hats are outer variable terms labelled by the entire terms \(p,q,r\).
Consequently,

\[
U=\widehat p+\bigl(\widehat q\cdot\widehat r\bigr)
\in T_\Sigma(T_\Sigma X).
\]

The recursive clauses calculate its value one constructor at a time:

\[
\begin{aligned}
\mu_X(U)
&=\mu_X(\widehat p)
  +\bigl(\mu_X(\widehat q)\cdot\mu_X(\widehat r)\bigr)\\
&=p+(q\cdot r)\\
&=(x+1)+\bigl(S(y)\cdot(x\cdot y)\bigr).
\end{aligned}
\]

The first recursive clause removes each hat; the second rebuilds the displayed
sum and product around the resulting inner terms. Thus ``flatten'' names the
fully specified recursive operation above.

The unit law \(\mu_X\circ\eta_{T_\Sigma X}=1_{T_\Sigma X}\) removes an outer
variable labelled by a whole term. The other unit law
\(\mu_X\circ T_\Sigma(\eta_X)=1_{T_\Sigma X}\) replaces every variable by
its corresponding variable term. The associativity square says that a
three-layer term yields the same ordinary term whether the inner two layers or
the outer two layers are reduced first. Structural induction proves both unit
laws, associativity, and the naturality of \(\eta\) and \(\mu\). For example, if
\(A=\widehat p+\widehat q\) and \(B=\widehat r\) are elements of
\(T_\Sigma^2X\), then
\(W=\eta_{T_\Sigma^2X}(A)\cdot\eta_{T_\Sigma^2X}(B)\) belongs to
\(T_\Sigma^3X\). Put
\(\widehat{p+q}=\eta_{T_\Sigma X}(p+q)\). The two paths around the
associativity square are the explicit calculations

\[
\begin{aligned}
(\mu_X\circ T_\Sigma\mu_X)(W)
  &=\mu_X\bigl(\widehat{p+q}\cdot\widehat r\bigr)
   =(p+q)\cdot r,\\
(\mu_X\circ\mu_{T_\Sigma X})(W)
  &=\mu_X(A\cdot B)
   =(p+q)\cdot r.
\end{aligned}
\]

\section{Syntactic substitution as Kleisli
composition}\label{syntactic-substitution-as-kleisli-composition}

A simultaneous substitution from a variable set \(X\) into terms over a
variable set \(Y\) is a family \((s_x)_{x\in X}\), with
\(s_x\in T_\Sigma Y\). Such a family supplies exactly one term \(s_x\) for
every \(x\in X\); consequently, the rule \(x\mapsto s_x\) is a total,
single-valued function. Denote it by

\[
\sigma:X\longrightarrow T_\Sigma Y,
\qquad \sigma(x)=s_x.
\]

For the familiar finite notation, let \(V\) be a variable set, let
\(x_1,\ldots,x_n\in V\) be pairwise distinct, and choose
\(s_1,\ldots,s_n\in T_\Sigma V\). Define
\(\sigma=[x_1:=s_1,\ldots,x_n:=s_n]\) by

\[
\sigma(z)=
\begin{cases}
s_i,&z=x_i\text{ for the unique }i\in\{1,\ldots,n\},\\
\eta_V(z),&z\notin\{x_1,\ldots,x_n\}.
\end{cases}
\]

Pairwise distinctness makes the first clause single-valued. The two cases
exhaust \(V\), and every displayed value belongs to \(T_\Sigma V\). Hence the
display defines the function \(\sigma:V\to T_\Sigma V\). The substitution specified by the empty list is
\(\eta_V\); a one-element list denotes ordinary single-variable substitution.

The term \(\Sigma\)-algebra \(T_\Sigma X\) carries the formal operations
\(f_X(t_1,\ldots,t_n)=f(t_1,\ldots,t_n)\). A function between two term
algebras is a \(\Sigma\)-algebra homomorphism when it preserves all these
operations.

\paragraph{Unique extension theorem.}
Every simultaneous substitution \(\sigma:X\to T_\Sigma Y\) has a unique
\(\Sigma\)-algebra homomorphism
\(\sigma^*:T_\Sigma X\to T_\Sigma Y\) satisfying

\[
\begin{aligned}
\sigma^*(\eta_X(x))&=\sigma(x),\\
\sigma^*\bigl(f(t_1,\ldots,t_n)\bigr)
&=f\bigl(\sigma^*(t_1),\ldots,\sigma^*(t_n)\bigr).
\end{aligned}
\]

The first clause supplies the value on a variable term. The second supplies the
value on a composite term from the values already assigned to its immediate
subterms. Structural recursion proves existence. For uniqueness, let
\(h:T_\Sigma X\to T_\Sigma Y\) satisfy the same clauses. The variable clause
yields \(h(\eta_X(x))=\sigma^*(\eta_X(x))\). If equality holds on
\(t_1,\ldots,t_n\), preservation of \(f\) yields equality on
\(f(t_1,\ldots,t_n)\). Structural induction now proves \(h=\sigma^*\).

The extension property is typed by the triangle

\[
\begin{tikzcd}[row sep=large,column sep=large]
X \arrow[r,"\eta_X"] \arrow[dr,"\sigma"'] &
T_\Sigma X \arrow[d,dashed,"{\exists!\,\sigma^*}"] \\
& T_\Sigma Y.
\end{tikzcd}
\]

The dashed arrow is required to be a \(\Sigma\)-algebra homomorphism. Its
preservation of an \(n\)-ary symbol \(f\) is the commutativity of

\[
\begin{tikzcd}[row sep=large,column sep=large]
(T_\Sigma X)^n \arrow[r,"f_X"] \arrow[d,"(\sigma^*)^n"'] &
T_\Sigma X \arrow[d,"\sigma^*"] \\
(T_\Sigma Y)^n \arrow[r,"f_Y"'] & T_\Sigma Y.
\end{tikzcd}
\]

The unique extension is also given by the outer path in the commutative
triangle

\[
\begin{tikzcd}[row sep=large,column sep=large]
T_\Sigma X \arrow[r,"T_\Sigma\sigma"] \arrow[dr,"\sigma^*"'] &
T_\Sigma^2Y \arrow[d,"\mu_Y"] \\
& T_\Sigma Y.
\end{tikzcd}
\]

Indeed, \(\mu_Y\circ T_\Sigma\sigma\) sends a variable term
\(\eta_X(x)\) to \(\sigma(x)\), preserves every operation symbol, and hence
equals \(\sigma^*\) by uniqueness. Taking
\(X=T_\Sigma Y\) and \(\sigma=1_{T_\Sigma Y}\) yields
\(\mu_Y=(1_{T_\Sigma Y})^*\). Monad multiplication is therefore the unique
homomorphic extension that substitutes each term for the outer variable
carrying that term as its label.

\paragraph{Example: a simultaneous substitution.}
Let \(X=\{x,y\}\), \(Y=\{u,v\}\), and define the function
\(\sigma:X\to T_\Sigma Y\) by
\(\sigma(x)=u+1\) and \(\sigma(y)=u\cdot v\). For
\(t=x\cdot(y+1)\), the recursive clauses yield

\[
\sigma^*(t)=(u+1)\cdot\bigl((u\cdot v)+1\bigr).
\]

The factorization through a term of terms is visible before multiplication:

\[
T_\Sigma\sigma(t)
=\eta_{T_\Sigma Y}(u+1)\cdot
 \bigl(\eta_{T_\Sigma Y}(u\cdot v)+1\bigr)
\in T_\Sigma(T_\Sigma Y).
\]

Applying \(\mu_Y\) replaces the two displayed outer variables by their term
labels. The word ``simultaneous'' means that the terms \(\sigma(x)\) and
\(\sigma(y)\) are inserted in one operation. For example,
\([x:=y,y:=S(y)]^*(x+y)=y+S(y)\), whereas first replacing \(x\) by \(y\)
and then replacing \(y\) by \(S(y)\) yields \(S(y)+S(y)\).

If \(\sigma:X\to T_\Sigma Y\) and \(\tau:Y\to T_\Sigma Z\), then their
Kleisli composite is the outer path in

\[
\begin{tikzcd}[column sep=large]
X \arrow[r,"\sigma"]
  \arrow[rrr,bend right=18,"\tau\star\sigma"'] &
T_\Sigma Y \arrow[r,"T_\Sigma\tau"] &
T_\Sigma^2Z \arrow[r,"\mu_Z"] &
T_\Sigma Z.
\end{tikzcd}
\]

Let \(\upsilon:Z\to T_\Sigma W\). The associativity of iterated syntactic
substitution is associativity of composition in the Kleisli category: the following
square commutes in \(\mathbf{Kl}(T_\Sigma)\), where each displayed arrow is a
Kleisli arrow:

\[
\begin{tikzcd}[row sep=large,column sep=huge]
X \arrow[r,"\tau\star\sigma"] \arrow[d,"\sigma"'] &
Z \arrow[d,"\upsilon"] \\
Y \arrow[r,"\upsilon\star\tau"'] & W.
\end{tikzcd}
\]

The upper-right path is \(\upsilon\star(\tau\star\sigma)\), and the
lower-left path is \((\upsilon\star\tau)\star\sigma\).

Ordinary replacement of a single variable is a special case. Given
\(s\in T_\Sigma X\) and \(x\in X\), define
\([x:=s]:X\to T_\Sigma X\) by

\[
[x:=s](y)=
\begin{cases}
s,&y=x,\\
\eta_X(y),&y\ne x.
\end{cases}
\]

Then \(t[x:=s]\) is the value of the Kleisli extension \([x:=s]^*\) at
\(t\).

\section{Structures as algebras for the term
monad}\label{structures-as-algebras-for-the-term-monad}

A \(\Sigma\)-algebra on a set \(A\) assigns an operation
\(f^{\mathcal A}:A^n\to A\) to every \(n\)-ary function symbol
\(f\in\Sigma\). Equivalently, it is an Eilenberg--Moore algebra for the
term monad: a map \(a:T_\Sigma A\to A\) satisfying the unit triangle and
multiplication square displayed in Section~\ref{categorical-vocabulary}, with
\(T=T_\Sigma\).

The map \(a\) evaluates a term whose variables have already been
replaced by elements of \(A\). The unit law says that an element used as
a variable evaluates to itself. The multiplication law says that
evaluating a term of terms in one step agrees with first evaluating the
inner terms and then the outer term.

For a language containing relation symbols, this algebra is only the
functional reduct of an \(\mathcal L\)-structure, obtained by forgetting
the interpretations of the relation symbols.

\Needspace{9\baselineskip}
An
\(\mathcal L\)-structure \(\mathcal A\) consists of:

\begin{itemize}
\tightlist
\item
  a nonempty set \(A\);
\item
  a \(T_\Sigma\)-algebra \(a:T_\Sigma A\to A\); and
\item
  a subset \(R^{\mathcal A}\subseteq A^n\) for every \(n\)-ary relation
  symbol \(R\in\mathcal R\).
\end{itemize}

Logical equality is interpreted by the diagonal
\(\Delta_A\subseteq A\times A\).

\section{Assignments and
evaluation}\label{assignments-and-evaluation}

Let \(X\) be a set of variables. An assignment in \(\mathcal A\) is a map
\(\alpha:X\to A\).

Term evaluation under \(\alpha\) is the outer path in the diagram

\[
\begin{tikzcd}[column sep=huge]
T_\Sigma X \arrow[r,"T_\Sigma\alpha"]
  \arrow[rr,bend right=18,"\llbracket-\rrbracket^{\mathcal A}_\alpha"'] &
T_\Sigma A \arrow[r,"a"] & A.
\end{tikzcd}
\]

The evaluation map \(\llbracket-\rrbracket^{\mathcal A}_\alpha\) is a
homomorphism from the term Eilenberg--Moore algebra
\((T_\Sigma X,\mu_X)\) to \((A,a)\); that assertion has the typed form

\[
\begin{tikzcd}[row sep=large,column sep=large]
T_\Sigma^2X
  \arrow[r,"T_\Sigma(\llbracket-\rrbracket^{\mathcal A}_\alpha)"]
  \arrow[d,"\mu_X"'] &
T_\Sigma A \arrow[d,"a"] \\
T_\Sigma X
  \arrow[r,"\llbracket-\rrbracket^{\mathcal A}_\alpha"'] & A.
\end{tikzcd}
\]

The evaluator extends the assignment: the triangle

\[
\begin{tikzcd}[row sep=large,column sep=large]
X \arrow[r,"\eta_X"] \arrow[dr,"\alpha"'] &
T_\Sigma X \arrow[d,"\llbracket-\rrbracket^{\mathcal A}_\alpha"] \\
& A
\end{tikzcd}
\]

commutes. Thus \(\llbracket x\rrbracket^{\mathcal A}_\alpha=\alpha(x)\), and

\[
\llbracket f(t_1,\ldots,t_n)\rrbracket^{\mathcal A}_\alpha
=
f^{\mathcal A}
\bigl(
\llbracket t_1\rrbracket^{\mathcal A}_\alpha,
\ldots,
\llbracket t_n\rrbracket^{\mathcal A}_\alpha
\bigr).
\]

The same structural-recursion and structural-induction argument used for
\(\sigma^*\) shows that this evaluator is the unique
\(\Sigma\)-algebra homomorphism extending \(\alpha\).
Because this extension property holds for every \(\Sigma\)-algebra
\((A,a)\) and every assignment \(\alpha:X\to A\),
\((T_\Sigma X,\mu_X)\) is called the \textbf{free} \(\Sigma\)-algebra on
\(X\).

For calculations, write the assignment as \(\atsign:X\to A\) and put
\(\atsign\,t:=\llbracket t\rrbracket^{\mathcal A}_{\atsign}\). Thus
\(\atsign\,x=\atsign(x)\) and
\(\atsign\,f(t_1,\ldots,t_n)=
f^{\mathcal A}(\atsign\,t_1,\ldots,\atsign\,t_n)\). When the structure must
be shown, write \(\atsign_{\mathcal A}\,t\).

Now let \(\sigma:X\to T_\Sigma Y\) be a syntactic substitution and let
\(\beta:Y\to A\) be an assignment. Substitution followed by evaluation
induces the assignment \(\beta\bullet\sigma:X\to A\), shown as the outer
path in

\[
\begin{tikzcd}[column sep=large]
X \arrow[r,"\sigma"]
  \arrow[rrr,bend right=18,"\beta\mathbin\bullet\sigma"'] &
T_\Sigma Y \arrow[r,"T_\Sigma\beta"] &
T_\Sigma A \arrow[r,"a"] & A.
\end{tikzcd}
\]

For each \(x\in X\), the value \((\beta\bullet\sigma)(x)\) is the value
under \(\beta\) of the term \(\sigma(x)\).

\paragraph{Example: evaluating a simultaneous substitution.}
Continue the earlier example with \(X=\{x,y\}\), \(Y=\{u,v\}\),
\(\sigma(x)=u+1\), \(\sigma(y)=u\cdot v\), and
\(t=x\cdot(y+1)\). In the arithmetic structure on \(\mathbb N\), let
\(\atsign:Y\to\mathbb N\) satisfy \(\atsign(u)=2\) and \(\atsign(v)=3\).
Then \((\atsign\bullet\sigma)(x)=3\) and
\((\atsign\bullet\sigma)(y)=6\), and

\[
\atsign\bigl(\sigma^*(t)\bigr)
=3\cdot(6+1)=21
= (\atsign\bullet\sigma)t.
\]

The term-substitution lemma says
that the square

\[
\begin{tikzcd}[row sep=large,column sep=huge]
T_\Sigma X
  \arrow[r,"\sigma^*"]
  \arrow[d,"\llbracket-\rrbracket^{\mathcal A}_{\beta\bullet\sigma}"'] &
T_\Sigma Y
  \arrow[d,"\llbracket-\rrbracket^{\mathcal A}_{\beta}"] \\
A \arrow[r,"1_A"'] & A
\end{tikzcd}
\]

commutes. Equivalently,
\(\llbracket\sigma^*(t)\rrbracket^{\mathcal A}_\beta=
\llbracket t\rrbracket^{\mathcal A}_{\beta\bullet\sigma}\).
\Needspace{11\baselineskip}
It follows from naturality of \(\mu\) and the algebra identity represented by
the two adjacent commutative squares:

\[
\begin{tikzcd}[row sep=large,column sep=large]
T_\Sigma X \arrow[r,"T_\Sigma\sigma"] &
T_\Sigma^2Y
  \arrow[r,"T_\Sigma^2\beta"]
  \arrow[d,"\mu_Y"'] &
T_\Sigma^2A
  \arrow[r,"T_\Sigma a"]
  \arrow[d,"\mu_A"'] &
T_\Sigma A \arrow[d,"a"] \\
& T_\Sigma Y
  \arrow[r,"T_\Sigma\beta"'] &
T_\Sigma A \arrow[r,"a"'] & A.
\end{tikzcd}
\]

For the single-variable substitution \([x:=s]\), the induced assignment is
\(\alpha\bullet[x:=s]=
\alpha[x\mapsto\llbracket s\rrbracket^{\mathcal A}_\alpha]\).

Consequently,

\[
\boxed{
\llbracket t[x:=s]\rrbracket^{\mathcal A}_\alpha
=
\llbracket t\rrbracket^{\mathcal A}_{
\alpha[x\mapsto\llbracket s\rrbracket^{\mathcal A}_\alpha]}
}
\]

For \(c\in A\), let \(\atsign\vert_c^x\) be the assignment obtained from
\(\atsign\) by sending \(x\) to \(c\) and leaving every other value fixed.
For a term \(s\), write \(\vert_s^x t\) for \(t[x:=s]\). The same substitution law is then the
type-revealing calculational identity

\[
\boxed{
\bigl(\atsign\vert_{\atsign\,s}^x\bigr)t
=\atsign(\vert_s^x t).
}
\]

The commutative square records the types and compositional structure of the
law; the \(\atsign\)-identity performs the corresponding calculation.
Syntactic substitution is Kleisli composition, and its compatibility with
evaluation is an Eilenberg--Moore algebra law. A concrete arithmetic comparison
of the textbook, categorical, and \(\atsign\)-forms appears below.

\section{Formulas and satisfaction}\label{formulas-and-satisfaction}

Let \(\operatorname{Form}_{\mathcal L}(X)\) be the set of
\(\alpha\)-equivalence classes of \(\mathcal L\)-formulas whose free
variables belong to \(X\). Two formulas are \(\alpha\)-equivalent when they
differ only by capture-avoiding renaming of bound variables. A Kleisli substitution
\(\sigma:X\to T_\Sigma Y\) induces capture-avoiding substitution
\((-)[\sigma]:\operatorname{Form}_{\mathcal L}(X)
\to\operatorname{Form}_{\mathcal L}(Y)\), which is well-defined on these
classes.

Put \(Q(X)=\operatorname{Form}_{\mathcal L}(X)\) and
\(Q(\sigma)=(-)[\sigma]\). On atomic formulas, \(Q(\sigma)\) substitutes
into their constituent terms. It commutes with the propositional connectives.
At a quantifier, bound variables are renamed when required to prevent capture.
The identity law is the commutativity of

\[
\begin{tikzcd}[row sep=large,column sep=large]
& Q(X) \arrow[dl,"Q(\eta_X)"'] \arrow[dr,"1_{Q(X)}"] & \\
Q(X) \arrow[rr,"1_{Q(X)}"'] && Q(X),
\end{tikzcd}
\]

and, for \(\sigma:X\to T_\Sigma Y\) and
\(\tau:Y\to T_\Sigma Z\), the composition law is the commutativity of

\[
\begin{tikzcd}[row sep=large,column sep=huge]
Q(X) \arrow[r,"Q(\sigma)"] \arrow[dr,"Q(\tau\star\sigma)"'] &
Q(Y) \arrow[d,"Q(\tau)"] \\
& Q(Z).
\end{tikzcd}
\]

Thus formulas carry an action of the Kleisli substitutions of the term
monad: \(Q:\mathbf{Kl}(T_\Sigma)\to\mathbf{Set}\) is a functor. Bound
variables add the capture-avoidance requirement to this action on formulas;
its composition laws remain those displayed above.

Satisfaction is defined in the usual way. Its atomic clauses are

\[
\begin{aligned}
\mathcal A,\alpha\models t=u
&\quad\Longleftrightarrow\quad
\llbracket t\rrbracket^{\mathcal A}_\alpha
=\llbracket u\rrbracket^{\mathcal A}_\alpha,\\
\mathcal A,\alpha\models R(t_1,\ldots,t_n)
&\quad\Longleftrightarrow\quad
\bigl(\llbracket t_1\rrbracket^{\mathcal A}_\alpha,\ldots,
\llbracket t_n\rrbracket^{\mathcal A}_\alpha\bigr)
\in R^{\mathcal A}.
\end{aligned}
\]

For calculations, define the postfix notation
\(\varphi\,\atsign:\!\Longleftrightarrow
\mathcal A,\atsign\models\varphi\). The atomic clauses become

\[
(t=u)\,\atsign\ \Longleftrightarrow\ \atsign\,t=\atsign\,u,
\qquad
R(t_1,\ldots,t_n)\,\atsign\ \Longleftrightarrow\ 
(\atsign\,t_1,\ldots,\atsign\,t_n)\in R^{\mathcal A}.
\]

When the structure must be shown, write
\(\varphi\,\atsign_{\mathcal A}\).

The clauses for the propositional connectives use the corresponding
operations in the metalanguage. Let \(\alpha:X\to A\). For \(b\in A\), write
\(\alpha[x\mapsto b]:X\cup\{x\}\to A\) for the assignment that sends
\(x\) to \(b\) and agrees with \(\alpha\) elsewhere; when \(x\in X\), this
overwrites its old value. If the free variables of \(\varphi\) lie in
\(X\cup\{x\}\), the quantifier clauses are

\[
\begin{aligned}
\mathcal A,\alpha\models(\forall x)\varphi
&\quad\Longleftrightarrow\quad
\text{for every }b\in A,
\ \mathcal A,\alpha[x\mapsto b]\models\varphi,\\
\mathcal A,\alpha\models(\exists x)\varphi
&\quad\Longleftrightarrow\quad
\text{for some }b\in A,
\ \mathcal A,\alpha[x\mapsto b]\models\varphi.
\end{aligned}
\]

Using \(\atsign\vert_b^x:X\cup\{x\}\to A\) for the assignment that sends
\(x\) to \(b\) and otherwise agrees with \(\atsign\), the same clauses are

\[
\begin{aligned}
((\forall x)\varphi)\,\atsign
&\quad\Longleftrightarrow\quad
\text{for every }b\in A,\ \varphi\,(\atsign\vert_b^x),\\
((\exists x)\varphi)\,\atsign
&\quad\Longleftrightarrow\quad
\text{for some }b\in A,\ \varphi\,(\atsign\vert_b^x).
\end{aligned}
\]

Regard satisfaction as its characteristic map
\(S_X^{\mathcal A}:\operatorname{Form}_{\mathcal L}(X)\times A^X
\to\mathbf 2\), with value \(\top\) exactly on satisfied pairs. For an
assignment \(\alpha:X\to A\), write
\(S_\alpha^{\mathcal A}:=S_X^{\mathcal A}(-,\alpha)\). For
\(\sigma:X\to T_\Sigma Y\) and \(\beta:Y\to A\), the formula-substitution
lemma is the commutativity of

\[
\begin{tikzcd}[row sep=large,column sep=large]
\operatorname{Form}_{\mathcal L}(X)
  \arrow[r,"{(-)[\sigma]}"]
  \arrow[d,swap,"S_{\beta\bullet\sigma}^{\mathcal A}"] &
\operatorname{Form}_{\mathcal L}(Y)
  \arrow[d,"S_\beta^{\mathcal A}"] \\
\mathbf 2 \arrow[r,"1_{\mathbf 2}"'] & \mathbf 2.
\end{tikzcd}
\]

Elementwise, the square says

\[
\boxed{
\mathcal A,\beta\models\varphi[\sigma]
\quad\Longleftrightarrow\quad
\mathcal A,\beta\bullet\sigma\models\varphi
}
\]

provided substitution is capture-avoiding. Its atomic cases are
consequences of the algebra law for term evaluation; its quantifier
cases follow by induction after the required renaming of bound
variables.

For a single-variable substitution, write \(\varphi\vert_s^x\) for the
capture-avoiding formula \(\varphi[x:=s]\). With the assignment-update notation
introduced above, the same lemma is the calculational identity

\[
\boxed{
\bigl(\varphi\vert_s^x\bigr)\,\atsign
\ \Longleftrightarrow\ 
\varphi\,\bigl(\atsign\vert_{\atsign\,s}^x\bigr).
}
\]

\paragraph{Example: avoiding capture.}
In arithmetic with \(<\), let \(\varphi\) be \((\exists y)\,x<y\), with
\(\operatorname{FV}(\varphi)=\{x\}\), and substitute the variable \(y\) for \(x\).
Capture-avoiding substitution yields
\(\varphi[x:=y]=(\exists z)\,y<z\), with \(z\) fresh. Textual
replacement without renaming the bound variable would yield \((\exists y)\,y<y\), 
changing the meaning. If \(\atsign(y)=3\) in \(\mathbb N\), then
\((\varphi\vert_y^x)\,\atsign\Longleftrightarrow
\varphi\,(\atsign\vert_{\atsign\,y}^x)\), and both sides are true; the formula produced
by replacement without renaming is false.

Satisfaction also has the locality property. If assignments
\(\alpha,\alpha':X\to A\) agree on every free variable of \(\varphi\),
then \(\mathcal A,\alpha\models\varphi\) if and only if
\(\mathcal A,\alpha'\models\varphi\).

In textbook usage, the notation \(\varphi(x_1,\ldots,x_n)\), when introduced
by saying that the free variables of \(\varphi\) are \emph{among}
\(x_1,\ldots,x_n\), means
\(\operatorname{FV}(\varphi)\subseteq\{x_1,\ldots,x_n\}\).

Some, all, or none of the displayed variables may actually occur freely.

\section{The at-sign notation as a pairing calculus}
\label{at-sign-pairing-calculus}

The assignment, term-evaluation, and satisfaction maps defined above have two
parallel presentations. Category-theoretic notation records their types and
composition; the $\atsign$-notation is used for their calculations. Prefix or
postfix position identifies the typed pairing in use.

\Needspace{9\baselineskip}
In this note, a \textbf{pairing} means simply a map $P\times Q\to K$ with two
ordered inputs. The two pairings are

\[
E_X^{\mathcal A}:A^X\times T_\Sigma X\longrightarrow A,
\qquad
S_X^{\mathcal A}:
\operatorname{Form}_{\mathcal L}(X)\times A^X
\longrightarrow\mathbf 2.
\]

Here $A^X$ is the set of assignments $X\to A$,
$E_X^{\mathcal A}(\alpha,t)=\llbracket t\rrbracket^{\mathcal A}_\alpha$,
and, for $\mathbf 2=\{\bot,\top\}$,
$S_X^{\mathcal A}(\varphi,\alpha)=\top$ exactly when
$\mathcal A,\alpha\models\varphi$. Thus $S_X^{\mathcal A}$ is the
characteristic function of the ordinary satisfaction relation.

The prefix and postfix notation already used above abbreviates these pairings:
$\atsign\,t=E_X^{\mathcal A}(\atsign,t)$, while
$\varphi\,\atsign$ holds exactly when
$S_X^{\mathcal A}(\varphi,\atsign)=\top$.

The placement of $\atsign$ records the assignment's argument position and
therefore identifies the pairing being used. Prefix $\atsign$ invokes
$E_X^{\mathcal A}$ and returns an element of $A$; postfix $\atsign$ invokes
$S_X^{\mathcal A}$ by abbreviating the assertion
$S_X^{\mathcal A}(\varphi,\atsign)=\top$.

\paragraph{Example: the two result types.}
In the arithmetic structure on \(\mathbb N\), let \(\atsign(x)=2\) and
\(\atsign(y)=3\). Then \(\atsign\bigl(x\cdot(y+1)\bigr)=8\in\mathbb N\),
whereas \(\bigl(x\cdot(y+1)=8\bigr)\,\atsign\) is a true proposition.

The update and substitution symbols introduced above make substitution appear
to pass through an evaluation pairing, from the syntactic argument to the
assignment. For an arbitrary Kleisli substitution
$\sigma:X\to T_\Sigma Y$, write
$P_a(\sigma):A^Y\to A^X$ for the assignment map
$\beta\mapsto\beta\bullet\sigma$. The compatibility of term evaluation with
substitution is the commutativity of the typed square

\[
\begin{tikzcd}[row sep=large,column sep=huge]
A^Y\times T_\Sigma X
  \arrow[r,"1_{A^Y}\times\sigma^*"]
  \arrow[d,"P_a(\sigma)\times1_{T_\Sigma X}"'] &
A^Y\times T_\Sigma Y \arrow[d,"E_Y^{\mathcal A}"] \\
A^X\times T_\Sigma X \arrow[r,"E_X^{\mathcal A}"'] & A.
\end{tikzcd}
\]

The compatibility of satisfaction with substitution is the commutativity of

\[
\begin{tikzcd}[row sep=large,column sep=large]
\operatorname{Form}_{\mathcal L}(X)\times A^Y
  \arrow[r,"{(-)[\sigma]\times1_{A^Y}}"]
  \arrow[d,"1_{Q(X)}\times P_a(\sigma)"'] &
\operatorname{Form}_{\mathcal L}(Y)\times A^Y
  \arrow[d,"S_Y^{\mathcal A}"] \\
\operatorname{Form}_{\mathcal L}(X)\times A^X
  \arrow[r,"S_X^{\mathcal A}"'] & \mathbf 2.
\end{tikzcd}
\]

In the $\atsign$-calculus, for \(\atsign:Y\to A\), the same two laws are

\[
(\atsign\bullet\sigma)t
=\atsign\bigl(\sigma^*(t)\bigr),
\qquad
\bigl(\varphi[\sigma]\bigr)\,\atsign
\ \Longleftrightarrow\ 
\varphi\,(\atsign\bullet\sigma).
\]

The lower-right object in the satisfaction square is \(\mathbf 2\), so its
$\atsign$-form is a biconditional. The diagrams and the $\atsign$-display
express the same compatibility of substitution with the evaluation pairings:
the diagrams expose every type, while the $\atsign$-forms perform the
corresponding calculations.

\subsection{One calculation in three notations}

Work in \(\mathbb N\) with \(X=\{x,y\}\), \(t=x\cdot(y+1)\),
\(s=y+1\), \(\alpha(x)=2\), and \(\alpha(y)=3\). Put
\(\sigma=[x:=s]\). Since \(\alpha\bullet\sigma=\alpha[x\mapsto4]\), the
same calculation has three forms:

\begin{itemize}
\tightlist
\item
  In traditional logical notation,
  \(\llbracket t[x:=s]\rrbracket^{\mathbb N}_\alpha=
  \llbracket t\rrbracket^{\mathbb N}_{\alpha[x\mapsto4]}=16\).
\item
  In the $\atsign$-calculus, the same equality is
  \(\bigl(\atsign\vert_{\atsign\,s}^x\bigr)t
  =\atsign(\vert_s^x t)=16\), where \(\atsign=\alpha\).
\item
  Categorically, the substitution square says
  \(E_X^{\mathbb N}(\alpha,\sigma^*(t))=
  E_X^{\mathbb N}(\alpha\bullet\sigma,t)\): \(\sigma^*\) changes the syntax,
  while \(\alpha\bullet\sigma\) changes the assignment.
\end{itemize}

Now take \(\varphi=(\exists z)\,x<z\), with
\(\operatorname{FV}(\varphi)=\{x\}\). In traditional notation,
\(\mathbb N,\alpha\models\varphi[x:=s]\) if and only if
\(\mathbb N,\alpha\bullet\sigma\models\varphi\), where
\(\alpha\bullet\sigma=\alpha[x\mapsto4]\). The categorical line is
\(S_X^{\mathbb N}(\varphi[\sigma],\alpha)=
S_X^{\mathbb N}(\varphi,\alpha\bullet\sigma)=\top\), and the pairing line is
\(\bigl(\varphi\vert_s^x\bigr)\,\atsign\Longleftrightarrow
\varphi\,\bigl(\atsign\vert_{\atsign\,s}^x\bigr)\). The categorical line identifies the common
substitution mechanism, while the $\atsign$-line performs the corresponding
calculation. Both remain alongside the textbook notation.

\section{Occurrence-indexed substitution}
\label{occurrence-indexed-substitution}

Term and formula substitution construct a new syntactic expression. A second
operation records where each occurrence of a formal argument goes during that
construction. \DP\ use finite ordinals for this bookkeeping
\cite[\S2.5, pp.~45--46]{DPcoherence}. Their operation is especially informative
when an outer pattern repeats an argument or when an inserted pattern repeats
one of its own arguments. Appendix~\ref{app:occurrence-feferman}
applies the same finite occurrence construction to arithmetic term holes
and proves its compatibility with Feferman's substitution on codes.

The letters \(\mu\) and \(\nu_i\) denote occurrence maps in the original
\DP\ notation. This note uses \(\mu_X\) for monad multiplication and reserves
\(\nu_A\) for the naming map introduced later. Accordingly, write \(o_M\) and
\(o_i\) for the occurrence maps in this section.

\subsection{Finite occurrence blocks}

For \(n\in\mathbb N\), let
\(\finord n=\{1,\ldots,n\}\), with
\(\finord0=\varnothing\). Given a finite list
\(\vec a=(a_1,\ldots,a_m)\), define its tagged block sum by

\[
\bigoplus_{r=1}^{m}\finord{a_r}
=\{(r,t):r\in\finord m,\ t\in\finord{a_r}\}.
\]

This tagged disjoint union is the coproduct in \(\mathbf{Set}\) of the finite
blocks \(\finord{a_r}\).
The tag \(r\) records the block and \(t\) records the position inside that
block. Put \(|\vec a|=\sum_{r=1}^{m}a_r\). The block encoding

\[
\theta_{\vec a}:
\bigoplus_{r=1}^{m}\finord{a_r}
\longrightarrow \finord{|\vec a|},
\qquad
\theta_{\vec a}(r,t)=\sum_{q<r}a_q+t
\]

is a bijection. For \(x\in\finord{|\vec a|}\), there is a unique block
\(r\) satisfying
\(\sum_{q<r}a_q<x\leq\sum_{q\leq r}a_q\); its inverse is

\[
\theta_{\vec a}^{-1}(x)
=\left(r,x-\sum_{q<r}a_q\right).
\]

For example, when \(\vec a=(2,1,3)\), the successive pairs
\((1,1),(1,2),(2,1),(3,1),(3,2),(3,3)\) are encoded by
\(1,2,3,4,5,6\). Finite-ordinal treatments call \(\theta_{\vec a}\)
\emph{flattening}. Its source is a tagged sum of finite sets. The term-monad
map \(\mu_X\) has source \(T_\Sigma(T_\Sigma X)\) and removes a syntactic term
layer. The two types keep the operations distinct.

The common zero-based convention
\([0]=\varnothing\) and \([n]=\{0,\ldots,n-1\}\) for \(n\geq 1\)
yields the same construction. The order isomorphism
\(\rho_n:\finord n\to[n]\), \(\rho_n(t)=t-1\), simply subtracts one from
each coordinate. In particular,
\(\rho_{|\vec a|}(\theta_{\vec a}(r,t))
=\sum_{q<r}a_q+(t-1)\), which is the zero-based block coordinate of the
local position \(t-1\in[a_r]\). Mac Lane uses this zero-based convention,
writing \(n\) for the finite ordinal denoted here by \([n]\)
\cite[\S VII.5]{MacLane}.

\subsection{Occurrence profiles and fibrewise substitution}

A \textbf{formula pattern} on a set \(P\) of placeholders is
an expression built from the members of \(P\), the truth constants
\(\top,\bot\), and propositional connectives. The truth constants are taken as
primitive nullary logical connectives or as fixed closed formula
abbreviations; in either case they determine elements of
\(\operatorname{Form}_{\mathcal L}(X)\) for every variable set \(X\). Write
\(\operatorname{Pat}(P)\) for the set of such patterns. Replacing its
placeholders by patterns or formulas while preserving the truth constants and
connectives produces another pattern or a formula. If
\(P=\{p_1,\ldots,p_k\}\), the pattern's \textbf{formal arity} is \(k\). Its
\textbf{occurrence slots} are the individual appearances of those
placeholders, counted with repetitions and read from left to right.

Let \(M\in\operatorname{Pat}(P)\) occupy \(m\) occurrence slots. Its
occurrence profile is the function

\[
o_M:\finord m\longrightarrow\finord k,
\qquad
o_M(r)=i
\quad\Longleftrightarrow\quad
\text{the \(r\)-th occurrence is occupied by \(p_i\)}.
\]

For comparison with the original \DP\ notation, let
\(\operatorname{sk}_P:\operatorname{Pat}(P)\to
\operatorname{Pat}(\{\square\})\) be the erasure map defined recursively by
replacing every formal placeholder in \(P\) by the single placeholder
\(\square\), while fixing truth constants and preserving connectives. When
the placeholder set is clear, write \(\operatorname{sk}\). The decorated
skeleton \(\operatorname{sk}_P(M)^{o_M}\) consists of the one-placeholder
skeleton together with the external argument-indexing map \(o_M\). Thus
\(M\) and \(\operatorname{sk}_P(M)^{o_M}\) encode the same pattern data in
two forms.

For each \(i\in\finord k\), let
\(Q_i=\{q_{i,1},\ldots,q_{i,\ell_i}\}\), with all these placeholders pairwise
distinct, and let \(N_i\in\operatorname{Pat}(Q_i)\) have \(n_i\) occurrence
slots. Its occurrence profile is a function
\(o_i:\finord{n_i}\to\finord{\ell_i}\). Put
\(Q_{\mathrm{occ}}=\coprod_{i=1}^k Q_i\), identified with the tagged set of placeholders
\(\{(i,q_{i,s}):i\in\finord k,\ s\in\finord{\ell_i}\}\).

Define \(\Phi_0:P\to\operatorname{Pat}(Q_{\mathrm{occ}})\) by
\(\Phi_0(p_i)=N_i\), viewing \(N_i\) in
\(\operatorname{Pat}(Q_{\mathrm{occ}})\) through the
inclusion \(Q_i\to Q_{\mathrm{occ}}\). Structural recursion yields a unique map
\((-)[\Phi_0]:\operatorname{Pat}(P)\to
\operatorname{Pat}(Q_{\mathrm{occ}})\) that sends each
\(p_i\) to \(N_i\), fixes \(\top\) and \(\bot\), and preserves every
propositional connective. Write \(M[\Phi_0]\) for the resulting substituted
pattern. While forming \(M[\Phi_0]\), the \(r\)-th occurrence of \(M\) is
replaced by a copy of \(N_{o_M(r)}\).

Write \(\vec n=(n_1,\ldots,n_k)\),
\(\vec\ell=(\ell_1,\ldots,\ell_k)\), and
\(\vec o=(o_1,\ldots,o_k)\).

The occurrences in the substituted pattern form the tagged source

\[
D_{M,\vec n}
=\bigoplus_{r=1}^{m}\finord{n_{o_M(r)}}
=\{(r,t):r\in\finord m,\ t\in\finord{n_{o_M(r)}}\}.
\]

The pair \((r,t)\) means: first choose the \(r\)-th occurrence of \(M\), then
choose the \(t\)-th occurrence inside the inserted copy of
\(N_{o_M(r)}\). The formal arguments of all the inserted patterns form the
tagged target

\[
C_{\vec\ell}
=\bigoplus_{i=1}^{k}\finord{\ell_i}
=\{(i,s):i\in\finord k,\ s\in\finord{\ell_i}\}.
\]

The fibrewise occurrence-substitution function is

\[
\gamma_{M,\vec o}:D_{M,\vec n}\longrightarrow C_{\vec\ell},
\qquad
\gamma_{M,\vec o}(r,t)
=\bigl(o_M(r),o_{o_M(r)}(t)\bigr).
\]

This formula is well-defined. Membership
\(t\in\finord{n_{o_M(r)}}\) and the declared type of
\(o_{o_M(r)}\) imply
\(o_{o_M(r)}(t)\in\finord{\ell_{o_M(r)}}\); hence the resulting pair lies in
the \(o_M(r)\)-th target fibre.

Let
\(\iota_r:\finord{n_{o_M(r)}}\to D_{M,\vec n}\) and
\(j_i:\finord{\ell_i}\to C_{\vec\ell}\) be the functions
\(\iota_r(t)=(r,t)\) and \(j_i(s)=(i,s)\). They are the coproduct
injections of the tagged block sums. For each \(r\in\finord m\), the square

\[
\begin{tikzcd}[row sep=large,column sep=large]
\finord{n_{o_M(r)}} \arrow[r,"\iota_r"]
  \arrow[d,"o_{o_M(r)}"'] &
D_{M,\vec n} \arrow[d,"\gamma_{M,\vec o}"] \\
\finord{\ell_{o_M(r)}} \arrow[r,"j_{o_M(r)}"'] & C_{\vec\ell}
\end{tikzcd}
\]

commutes. Taken together for every \(r\), these squares determine
\(\gamma_{M,\vec o}\) uniquely, because every element of
\(D_{M,\vec n}\) lies in exactly one injected fibre.
This is the coproduct unique-extension property for the family of maps
\(j_{o_M(r)}\circ o_{o_M(r)}\).

Projection to the block tags yields another typed description. Define
\(p_D(r,t)=r\) and \(p_C(i,s)=i\). The following square commutes:

\[
\begin{tikzcd}[row sep=large,column sep=large]
D_{M,\vec n} \arrow[r,"\gamma_{M,\vec o}"] \arrow[d,"p_D"'] &
C_{\vec\ell} \arrow[d,"p_C"] \\
\finord m \arrow[r,"o_M"'] & \finord k
\end{tikzcd}
\]

On the fibre above \(r\), the map is exactly
\(o_{o_M(r)}\), with its value placed in the fibre above \(o_M(r)\).

\subsection{The ordinal-coded operation}

Define cumulative source and target sizes by

\[
\pi(0)=0,\qquad
\pi(r)=\sum_{q=1}^{r}n_{o_M(q)},
\qquad
\lambda(0)=0,\qquad
\lambda(i)=\sum_{q=1}^{i}\ell_q.
\]

Thus \(D_{M,\vec n}\) has \(\pi(m)\) elements and
\(C_{\vec\ell}\) has \(\lambda(k)\) elements. Their block encodings are

\[
\begin{aligned}
\theta_D:D_{M,\vec n}&\longrightarrow\finord{\pi(m)},
&\theta_D(r,t)&=\pi(r-1)+t,\\
\theta_C:C_{\vec\ell}&\longrightarrow\finord{\lambda(k)},
&\theta_C(i,s)&=\lambda(i-1)+s.
\end{aligned}
\]

The \DP\ ordinal-coded substitution is the transported function

\[
\delta_{M,\vec o}
=\theta_C\circ\gamma_{M,\vec o}\circ\theta_D^{-1}:
\finord{\pi(m)}\longrightarrow\finord{\lambda(k)}.
\]

Its relation to the fibrewise map is the commutative square

\[
\begin{tikzcd}[row sep=large,column sep=large]
D_{M,\vec n} \arrow[r,"\gamma_{M,\vec o}"] \arrow[d,"\theta_D"'] &
C_{\vec\ell} \arrow[d,"\theta_C"] \\
\finord{\pi(m)} \arrow[r,"\delta_{M,\vec o}"'] &
\finord{\lambda(k)}.
\end{tikzcd}
\]

The coordinate formula uses the block locator
\(b:\finord{\pi(m)}\to\finord m\), defined by

\[
b(x)=\min\{r\in\finord m:x\leq\pi(r)\}.
\]

For every \(x\in\finord{\pi(m)}\), the displayed set is nonempty because
\(x\leq\pi(m)\); its least member is therefore defined and unique.
The local position of \(x\) inside that source block is
\(x-\pi(b(x)-1)\). Consequently,

\[
\delta_{M,\vec o}(x)
=\lambda\bigl(o_M(b(x))-1\bigr)
+o_{o_M(b(x))}\bigl(x-\pi(b(x)-1)\bigr).
\]

To verify the formula, start with \(x\). The inverse block encoding satisfies
\(\theta_D^{-1}(x)=(b(x),x-\pi(b(x)-1))\). Applying \(\gamma_{M,\vec o}\)
selects the profile \(o_{o_M(b(x))}\). Applying \(\theta_C\) adds the target
offset \(\lambda(o_M(b(x))-1)\). This is exactly the displayed value of
\(\delta_{M,\vec o}(x)\).

\paragraph{Occurrence-profile theorem.}
Under the block encodings \(\theta_D\) and \(\theta_C\), the occurrence
profile of the substituted pattern \(M[\Phi_0]\) is
\(\delta_{M,\vec o}\). Consequently, the labelled-pattern construction
transports the \DP\ generalized composite to

\[
\operatorname{sk}_P(M)^{o_M}
\bigl(\operatorname{sk}_{Q_1}(N_1)^{o_1},\ldots,
      \operatorname{sk}_{Q_k}(N_k)^{o_k}\bigr)
:=\operatorname{sk}_{Q_{\mathrm{occ}}}\bigl(M[\Phi_0]\bigr)^{\delta_{M,\vec o}},
\]

where a superscript records the external argument-indexing map attached to the
displayed skeleton. Renaming the skeletons
\(\operatorname{sk}_P(M),\operatorname{sk}_{Q_1}(N_1),\ldots,
\operatorname{sk}_{Q_k}(N_k)\) as \(M,N_1,\ldots,N_k\), as in the original
notation, this is
\(M^\mu(N_1^{\nu_1},\ldots,N_k^{\nu_k})\), with
\(\mu=o_M\), \(\nu_i=o_i\), and
\(\mu(\nu_1,\ldots,\nu_k)=\delta_{M,\vec o}\).

Indeed, every occurrence of \(M[\Phi_0]\) has unique fibre coordinates
\((r,t)\in D_{M,\vec n}\): \(r\) selects an occurrence of \(M\), and \(t\)
selects an occurrence in the inserted copy of \(N_{o_M(r)}\). The placeholder
at that occurrence is the tagged element
\(\bigl(o_M(r),q_{o_M(r),\,o_{o_M(r)}(t)}\bigr)\in Q_{\mathrm{occ}}\).
Under the canonical indexing of \(Q_{\mathrm{occ}}\) by \(C_{\vec\ell}\), its
coordinate is \(\gamma_{M,\vec o}(r,t)\). Transporting source and target
coordinates through \(\theta_D\) and \(\theta_C\) yields
\(\theta_C\circ\gamma_{M,\vec o}\circ\theta_D^{-1}
=\delta_{M,\vec o}\), as required.

For maps \(f_i:\finord{a_i}\to\finord{b_i}\), put
\(\vec a=(a_1,\ldots,a_k)\) and \(\vec b=(b_1,\ldots,b_k)\). Their
flattened ordinal sum is

\[
f_1\oplus\cdots\oplus f_k
=\theta_{\vec b}\circ
 \left(\bigoplus_{i=1}^{k}f_i\right)
 \circ\theta_{\vec a}^{-1},
\qquad
\left(\bigoplus_{i=1}^{k}f_i\right)(i,t)=(i,f_i(t)).
\]

When \(m=k\) and \(o_M=1_{\finord k}\), each outer slot uses the inserted
pattern with the same index. The fibrewise map is
\(\bigoplus_{i=1}^{k}o_i\), and its ordinal-coded form is
\(o_1\oplus\cdots\oplus o_k\).

\subsection{A complete occurrence calculation}

Take the formula patterns

\[
M(p,q)=q\wedge(p\vee q),\qquad
N_1(a,b)=a\wedge b,
\qquad
N_2(c,d)=d\wedge(c\vee d).
\]

Reading occurrences from left to right yields
\(o_M=(2,1,2)\), \(o_1=(1,2)\), and \(o_2=(2,1,2)\).

Here \(m=3\), \(k=2\), \((n_1,n_2)=(2,3)\), and
\((\ell_1,\ell_2)=(2,2)\). After substitution,

\[
M(N_1,N_2)
=\bigl(d\wedge(c\vee d)\bigr)
 \wedge
 \bigl((a\wedge b)\vee(d\wedge(c\vee d))\bigr).
\]

Its source fibres have sizes \(3,2,3\), while its target fibres list the
formal arguments as \((a,b)\) and \((c,d)\). The successive source
occurrences and their encoded target positions are

\[
\begin{array}{c|ccc|cc|ccc}
\text{source position }x&1&2&3&4&5&6&7&8\\ \hline
\text{occurring letter}&d&c&d&a&b&d&c&d\\
\delta_{M,\vec o}(x)&4&3&4&1&2&4&3&4
\end{array}
\]

Thus the generalized composite in the labelled-pattern encoding of the
\DP\ notation is
\[
\operatorname{sk}_P(M)^{(2,1,2)}
\bigl(\operatorname{sk}_{Q_1}(N_1)^{(1,2)},
      \operatorname{sk}_{Q_2}(N_2)^{(2,1,2)}\bigr)
=\operatorname{sk}_{Q_{\mathrm{occ}}}
 \bigl(M(N_1,N_2)\bigr)^{(4,3,4,1,2,4,3,4)}.
\]

For instance, \(\theta_D^{-1}(2)=(1,2)\): source position \(2\) is the
second occurrence inside the copy of \(N_2\) inserted at the first outer
slot. Then
\(\gamma_{M,\vec o}(1,2)=(2,o_2(2))=(2,1)\), and
\(\theta_C(2,1)=3\). Hence \(\delta_{M,\vec o}(2)=3\), as the table records.

Three substitution operations now have explicit types.
The map \(\sigma^*:T_\Sigma X\to T_\Sigma Y\) substitutes terms for term
variables. The map
\((-)[\sigma]:\operatorname{Form}_{\mathcal L}(X)\to
\operatorname{Form}_{\mathcal L}(Y)\) performs capture-avoiding term
substitution in formulas. The map
\((-)[\Phi_0]:\operatorname{Pat}(P)\to
\operatorname{Pat}(Q_{\mathrm{occ}})\) substitutes
formula patterns for formal formula arguments. The maps
\(\gamma_{M,\vec o}\) and \(\delta_{M,\vec o}\) record the intermediate
pattern-level provenance of placeholder occurrences under the third operation.

\subsection{The \texorpdfstring{$\atsign$}{@}-calculus for the associated
formula substitution}

The \(\atsign\)-calculus evaluates the formulas produced by the pattern
substitution whose occurrence provenance is recorded by
\(\gamma_{M,\vec o}\) or \(\delta_{M,\vec o}\).
For every
\(i\in\finord k\) and \(s\in\finord{\ell_i}\), choose an actual
\(\mathcal L\)-formula \(\psi_{i,s}\) whose free variables belong to \(X\),
and define
\(\Psi_i:Q_i\to
\operatorname{Form}_{\mathcal L}(X)\) by
\(\Psi_i(q_{i,s})=\psi_{i,s}\). Recursively replace every
\(q_{i,s}\) in \(N_i\) by \(\psi_{i,s}\), fixing \(\top,\bot\), and
preserving each connective, and
write the resulting formula as \(\widehat N_i=N_i[\Psi_i]\). The outer
formula-pattern substitution is the function

\[
\Phi:P\longrightarrow\operatorname{Form}_{\mathcal L}(X),
\qquad \Phi(p_i)=\widehat N_i.
\]

Structural recursion yields a unique truth-constant-and-connective-preserving extension
\((-)[\Phi]:\operatorname{Pat}(P)\to
\operatorname{Form}_{\mathcal L}(X)\): it sends \(p_i\) to \(\Phi(p_i)\)
and sends each connective to the same connective applied to the recursively
substituted arguments. Write \(M[\Phi]\) for the image of \(M\) under this
extension.

The maps \(\Psi_i\) combine uniquely to
\(\Psi:Q_{\mathrm{occ}}\to\operatorname{Form}_{\mathcal L}(X)\).
Structural recursion yields its unique extension
\((-)[\Psi]:\operatorname{Pat}(Q_{\mathrm{occ}})\to
\operatorname{Form}_{\mathcal L}(X)\). Associativity of pattern substitution
is the map identity
\[
(-)[\Psi]\circ(-)[\Phi_0]=(-)[\Phi],
\]
and hence, at \(M\), the equation
\[
\bigl(M[\Phi_0]\bigr)[\Psi]=M[\Phi].
\]
Both sides are obtained by the unique extension that sends \(p_i\) to
\(N_i[\Psi_i]\), fixes \(\top,\bot\), and preserves every connective. The
left side is therefore the formula obtained by instantiating the pattern
underlying the \DP\ generalized composite.

For an \(\mathcal L\)-structure \(\mathcal A\) and an assignment
\(\atsign:X\to A\), define the induced truth assignment

\[
\atsign\mathbin\bullet\Phi:P\longrightarrow\mathbf 2,
\qquad
(\atsign\mathbin\bullet\Phi)(p_i)
=S_X^{\mathcal A}(\Phi(p_i),\atsign).
\]

Write \(\atsign_\Phi:=\atsign\mathbin\bullet\Phi\) for this induced truth
assignment.
This use of \(\mathbin\bullet\) is the truth-valued analogue of the induced
assignment \(\beta\mathbin\bullet\sigma\) for term substitution. The truth
evaluation pairing has type

\[
S_P^{\mathbf 2}:\operatorname{Pat}(P)\times\mathbf 2^P
\longrightarrow\mathbf 2.
\]

For \(v:P\to\mathbf2\), it is defined recursively by
\(S_P^{\mathbf2}(p_i,v)=v(p_i)\),
\(S_P^{\mathbf2}(\top,v)=\top\), and
\(S_P^{\mathbf2}(\bot,v)=\bot\). For example,
\(S_P^{\mathbf2}(\neg U,v)=\neg S_P^{\mathbf2}(U,v)\) and
\(S_P^{\mathbf2}(U\wedge V,v)=
S_P^{\mathbf2}(U,v)\wedge S_P^{\mathbf2}(V,v)\); every further connective
uses its corresponding truth function.

\Needspace{11\baselineskip}
For a truth assignment \(\atsign_P:P\to\mathbf 2\), let
\(S_{\atsign_P}^{\mathbf 2}=S_P^{\mathbf 2}(-,\atsign_P):
\operatorname{Pat}(P)\to\mathbf 2\), and write
\(M\,\atsign_P:\!\Longleftrightarrow
S_P^{\mathbf 2}(M,\atsign_P)=\top\). Substitution of
the formulas \(\Phi(p_i)\) for the formal arguments is compatible with the two
satisfaction pairings, as expressed by the following commutative square:

\[
\begin{tikzcd}[row sep=large,column sep=large]
\operatorname{Pat}(P) \arrow[r,"{(-)[\Phi]}"]
  \arrow[d,"S_{\atsign_\Phi}^{\mathbf 2}"'] &
\operatorname{Form}_{\mathcal L}(X)
  \arrow[d,"S_{\atsign}^{\mathcal A}"] \\
\mathbf 2 \arrow[r,"1_{\mathbf 2}"'] & \mathbf 2
\end{tikzcd}
\]

In the \(\atsign\)-calculus, this is the sliding law

\[
\boxed{
\bigl(M[\Phi]\bigr)\,\atsign
\quad\Longleftrightarrow\quad
M\,\atsign_\Phi.
}
\]

The law follows by structural induction on \(M\). At the placeholder
\(p_i\), it is the defining equation for
\(\atsign_\Phi=\atsign\mathbin\bullet\Phi\). At \(\top\) and \(\bot\), it is their
respective truth clause. Each connective step follows from the corresponding
truth function.

Combining this sliding law with pattern-substitution associativity applies the
\(\atsign\)-calculus directly to the instantiated \DP\ composite:
\[
\boxed{
\bigl((M[\Phi_0])[\Psi]\bigr)\,\atsign
\quad\Longleftrightarrow\quad
M\,\atsign_\Phi.
}
\]

The occurrence maps have a direct \(\atsign\)-labelled form. Give every
formal argument of every inserted pattern its truth value under the current
assignment by defining

\[
e_{\atsign}:C_{\vec\ell}\longrightarrow\mathbf 2,
\qquad
e_{\atsign}(i,s)=S_X^{\mathcal A}(\psi_{i,s},\atsign),
\qquad
d_{\atsign}:D_{M,\vec n}\longrightarrow\mathbf 2,
\quad d_{\atsign}=e_{\atsign}\circ\gamma_{M,\vec o}.
\]

In the postfix \(\atsign\)-notation, these definitions say
\[
\begin{aligned}
e_{\atsign}(i,s)=\top
&\quad\Longleftrightarrow\quad \psi_{i,s}\,\atsign,\\
d_{\atsign}(r,t)=\top
&\quad\Longleftrightarrow\quad
\psi_{o_M(r),\,o_{o_M(r)}(t)}\,\atsign.
\end{aligned}
\]

Thus \(d_{\atsign}\) labels each pattern-level placeholder occurrence in the
substituted pattern by pulling back the truth label of the formal argument from
which it came. Occurrences internal to an actual formula \(\psi_{i,s}\) are
represented here by the single truth value of that formula. In fibre
coordinates this is the commutative triangle

\[
\begin{tikzcd}[row sep=large,column sep=large]
D_{M,\vec n} \arrow[r,"\gamma_{M,\vec o}"]
  \arrow[dr,"d_{\atsign}"'] &
C_{\vec\ell} \arrow[d,"e_{\atsign}"] \\
& \mathbf 2.
\end{tikzcd}
\]

After the block encodings, put
\(\bar e_{\atsign}=e_{\atsign}\circ\theta_C^{-1}\) and
\(\bar d_{\atsign}=d_{\atsign}\circ\theta_D^{-1}\). The corresponding
ordinal-coordinate triangle is

\[
\begin{tikzcd}[row sep=large,column sep=large]
\finord{\pi(m)} \arrow[r,"\delta_{M,\vec o}"]
  \arrow[dr,"\bar d_{\atsign}"'] &
\finord{\lambda(k)} \arrow[d,"\bar e_{\atsign}"] \\
& \mathbf 2.
\end{tikzcd}
\qquad
\bar d_{\atsign}=\bar e_{\atsign}\circ\delta_{M,\vec o}.
\]

These diagrams are the direct interface between the \(\atsign\)-calculus and
the occurrence operation: truth evaluation supplies the pattern-level
formal-argument occurrence labels, and the
occurrence maps transport those labels through substitution. Evaluating the
connective tree of \(M[\Phi]\) then returns its single truth value.

For the preceding example, replace \(a,b,c,d\) by actual formulas
\(\psi_a,\psi_b,\psi_c,\psi_d\). Thus
\(\widehat N_1=\psi_a\wedge\psi_b\),
\(\widehat N_2=\psi_d\wedge(\psi_c\vee\psi_d)\),
\(\Phi(p)=\widehat N_1\), and \(\Phi(q)=\widehat N_2\). Suppose
\(\psi_a\,\atsign\), \(\psi_b\,\atsign\), and
\(\psi_c\,\atsign\) are true, while \(\psi_d\,\atsign\) is false. In the
target order \((a,b,c,d)\), the truth labels and their pullback along the
computed occurrence map are

\[
\bar e_{\atsign}=(\top,\top,\top,\bot),
\qquad
\bar d_{\atsign}
=\bar e_{\atsign}\circ\delta_{M,\vec o}
=(\bot,\top,\bot,\top,\top,\bot,\top,\bot).
\]

The induced assignment to the outer arguments satisfies
\(\atsign_\Phi(p)=\top\) and \(\atsign_\Phi(q)=\bot\). Hence

\[
S_P^{\mathbf 2}(M,\atsign_\Phi)
=\bot\wedge(\top\vee\bot)=\bot,
\]

and the sliding law says that \(\bigl(M[\Phi]\bigr)\,\atsign\) is false.


\section{Expansion of a language by names for a
structure}\label{expansion-of-a-language-by-names-for-a-structure}

Let \(\mathcal A\) be an \(\mathcal L\)-structure with domain \(A\). The
expansion of \(\mathcal L\) by names for elements of \(A\), written here
as \(\mathcal L_A\), adds a new constant symbol \(\bar c\) for every
\(c\in A\). The symbols \(\bar c\) form a disjoint copy of \(A\): the
element \(c\) lies in the domain, while its name \(\bar c\) is a formal
constant in the expanded signature.

Let \(\Sigma[A]\) be the expanded functional signature. A term over
\(\Sigma[A]\) with variables in \(X\) is equivalently an original
\(\Sigma\)-term whose leaves may be either named elements of \(A\) or
variables in \(X\).

\Needspace{8\baselineskip}
This description by leaves yields a natural bijection
\(\chi_X:T_{\Sigma[A]}X\xrightarrow{\cong}T_\Sigma(A+X)\).

Here \(A+X\) is the disjoint union, or coproduct in \(\mathbf{Set}\),
with canonical inclusions \(\operatorname{inl}:A\to A+X\) and
\(\operatorname{inr}:X\to A+X\). Under this identification, the
\(A\)-summand is held fixed by substitution, while the \(X\)-summand
contains the substitutable variables.

\Needspace{8\baselineskip}
This bijection is \textbf{natural in \(X\)}: for every \(u:X\to Y\), the
square

\[
\begin{tikzcd}[row sep=large,column sep=huge]
T_{\Sigma[A]}X \arrow[r,"\chi_X"] \arrow[d,"T_{\Sigma[A]}u"'] &
T_\Sigma(A+X) \arrow[d,"T_\Sigma(1_A+u)"] \\
T_{\Sigma[A]}Y \arrow[r,"\chi_Y"'] & T_\Sigma(A+Y)
\end{tikzcd}
\]

commutes. The monad structure on \(X\mapsto T_\Sigma(A+X)\) in this
description is the expanded term-monad structure transported through
\(\chi\).

The inclusion of the original signature into the expanded signature induces
the monad morphism \(j:T_\Sigma\Rightarrow T_{\Sigma[A]}\). Under the
displayed identification, its component is typed by the commutative triangle

\[
\begin{tikzcd}[row sep=large,column sep=huge]
T_\Sigma X \arrow[r,"j_X"]
  \arrow[dr,"T_\Sigma(\operatorname{inr})"'] &
T_{\Sigma[A]}X \arrow[d,"\chi_X"] \\
& T_\Sigma(A+X)
\end{tikzcd}
\]

Here \(\operatorname{inr}:X\to A+X\) includes the variable summand.
Thus \(j\) regards every original term as an expanded-language term
without inserting any new constants. Because \(j\) is a monad morphism,
this inclusion preserves variables and syntactic substitution.

\Needspace{9\baselineskip}
More generally, to give a \(\Sigma[A]\)-algebra structure on a set \(B\)
is to give:

\begin{itemize}
\tightlist
\item
  a \(T_\Sigma\)-algebra \(b:T_\Sigma B\to B\); and
\item
  a function \(k:A\to B\) interpreting the newly added constants.
\end{itemize}

The full expanded algebra map \(h_{b,k}:T_{\Sigma[A]}B\to B\) is the outer
path in

\[
\begin{tikzcd}[column sep=large]
T_{\Sigma[A]}B \arrow[r,"\chi_B"]
  \arrow[rrr,bend right=18,"{h_{b,k}}"'] &
T_\Sigma(A+B) \arrow[r,"{T_\Sigma[k,1_B]}"] &
T_\Sigma B \arrow[r,"b"] & B.
\end{tikzcd}
\]

Here \([k,1_B]:A+B\to B\) is the \textbf{copairing} of \(k\) and the
identity of \(B\): it sends the \(A\)-summand by \(k\) and the
\(B\)-summand by \(1_B\). It is the unique function with those two
restrictions.

The expanded algebra retains the original \(\Sigma\)-operations.

\Needspace{8\baselineskip}
Their preservation is expressed by the following commutative restriction
triangle:

\[
\begin{tikzcd}[row sep=large,column sep=huge]
T_\Sigma B \arrow[r,"j_B"] \arrow[dr,"b"'] &
T_{\Sigma[A]}B \arrow[d,"{h_{b,k}}"] \\
& B.
\end{tikzcd}
\]

The distinguished expansion \(\mathcal A^+\) of \(\mathcal A\) takes
\(B=A\), retains every original relation \(R^{\mathcal A}\), and uses the pair
\((a,1_A)\): the original algebra map \(a:T_\Sigma A\to A\) together with the
identity interpretation of the new names. Denote the resulting specialization
of the preceding path by \(a^+:T_{\Sigma[A]}A\to A\).
In particular, \(\bar c^{\mathcal A^+}=c\) for every \(c\in A\).

Restricting the functional reduct of \(\mathcal A^+\) along the monad
morphism \(j\) recovers the functional reduct of \(\mathcal A\):
this is the specialization \(B=A\), \(b=a\), and \(k=1_A\) of the restriction
triangle. Forgetting the newly named constants changes none of the original
operations.

\Needspace{7\baselineskip}
For ground terms, let \(\rho_A:A+\varnothing\xrightarrow{\cong}A\) be the
canonical bijection. Ground-term evaluation is the map

\[
\begin{tikzcd}[column sep=large]
T_{\Sigma[A]}\varnothing \arrow[r,"\chi_\varnothing"]
  \arrow[rrr,bend right=18,"e_A"'] &
T_\Sigma(A+\varnothing) \arrow[r,"T_\Sigma(\rho_A)"] &
T_\Sigma A \arrow[r,"a"] & A.
\end{tikzcd}
\]
After the canonical identification \(A+\varnothing\cong A\), this map is the
original algebra map \(a:T_\Sigma A\to A\).

The naming map \(\nu_A:A\to T_{\Sigma[A]}\varnothing\), given by
\(c\mapsto\bar c\), is related to the original monad unit by the commutative
square

\[
\begin{tikzcd}[row sep=large,column sep=huge]
A \arrow[r,"\nu_A"] \arrow[d,"\operatorname{inl}"'] &
T_{\Sigma[A]}\varnothing \arrow[d,"\chi_\varnothing"] \\
A+\varnothing \arrow[r,"\eta^{T_\Sigma}_{A+\varnothing}"'] &
T_\Sigma(A+\varnothing).
\end{tikzcd}
\]

\Needspace{10\baselineskip}
After the canonical identification \(A+\varnothing\cong A\), the square says
that \(\nu_A\) corresponds to \(\eta_A^{T_\Sigma}:A\to T_\Sigma A\). The
algebra unit triangle then makes the following section triangle commute:

\[
\begin{tikzcd}[row sep=large,column sep=huge]
A \arrow[r,"\nu_A"] \arrow[dr,"1_A"'] &
T_{\Sigma[A]}\varnothing \arrow[d,"e_A"] \\
& A.
\end{tikzcd}
\]

Thus every element has a chosen ground-term name, and evaluating that
name returns the element. A \textbf{ground term} is a term with no
variables. The section triangle says that \(\nu_A\) is a
\textbf{section}, or right inverse, of ground-term evaluation. It need
not be a two-sided inverse: many different ground terms can have the
same value in \(\mathcal A^+\).

Let \(\epsilon_A:\varnothing\to A\) be the unique empty assignment, and
write \(\atsign_{\varnothing,\mathcal A^+}:=\epsilon_A\) when evaluation
takes place in \(\mathcal A^+\). In this two-subscript notation for empty
assignments, the first subscript records the variable context and the second
specifies the structure. In the pairing calculus, the section equation says that
evaluating a name returns its value:
\(\atsign_{\varnothing,\mathcal A^+}\,\bar c=c\), where \(c\in A\).

\Needspace{5\baselineskip}
\section{Assignments as substitutions of
names}\label{assignments-as-substitutions-of-names}

The expansion \(\mathcal L_A\) converts an assignment into a syntactic
substitution by ground terms. Given \(\alpha:X\to A\), define
\(\bar\alpha:=\nu_A\circ\alpha:X\to T_{\Sigma[A]}\varnothing\).

\Needspace{7\baselineskip}
This definition is encoded by the commutative triangle

\[
\begin{tikzcd}[column sep=huge]
X \arrow[r,"\alpha"]
  \arrow[rr,bend right=18,"\bar\alpha"'] &
A \arrow[r,"\nu_A"] & T_{\Sigma[A]}\varnothing.
\end{tikzcd}
\]

The Kleisli substitution \(\bar\alpha\) sends each variable \(x\) to the
constant \(\overline{\alpha(x)}\). Regard every
original term as an expanded-language term by \(j_X\). The inclusion of
languages also induces
\(\iota_X:\operatorname{Form}_{\mathcal L}(X)\to
\operatorname{Form}_{\mathcal L_A}(X)\). Throughout this section, an original
formula is regarded as an expanded-language formula by \(\iota_X\). To keep the calculations readable,
write \(t[\bar\alpha]\) for \(\bar\alpha^*(j_X(t))\) and
\(\varphi[\bar\alpha]\) for
\(\iota_X(\varphi)[\bar\alpha]\).

\Needspace{14\baselineskip}
\paragraph{Categorical form.}
For \(t\in T_\Sigma X\) and
\(\varphi\in\operatorname{Form}_{\mathcal L}(X)\), evaluation factors through
the language inclusion and the Kleisli extension of \(\bar\alpha\). The two
factorizations are the outer paths in the following typed diagrams:

\[
\begin{tikzcd}[row sep=large,column sep=large]
T_\Sigma X
  \arrow[r,"j_X"]
  \arrow[drr,bend right=18,"{E_X^{\mathcal A}(\alpha,-)}"'] &
T_{\Sigma[A]}X \arrow[r,"\bar\alpha^*"] &
T_{\Sigma[A]}\varnothing \arrow[d,"e_A"] \\
&& A
\end{tikzcd}
\]

\[
\begin{tikzcd}[row sep=large,column sep=large]
\operatorname{Form}_{\mathcal L}(X)
  \arrow[r,"\iota_X"]
  \arrow[drr,bend right=18,"{S_X^{\mathcal A}(-,\alpha)}"'] &
\operatorname{Form}_{\mathcal L_A}(X)
  \arrow[r,"{(-)[\bar\alpha]}"] &
\operatorname{Form}_{\mathcal L_A}(\varnothing)
  \arrow[d,"{S_\varnothing^{\mathcal A^+}(-,\epsilon_A)}"] \\
&& \mathbf 2
\end{tikzcd}
\]

The first diagram applies ground-term evaluation to \(t[\bar\alpha]\). The
second evaluates the truth value of the sentence \(\varphi[\bar\alpha]\);
replacing every free variable by a name leaves the empty assignment as its only
assignment argument.

\paragraph{At-sign calculation.}
Set \(\atsign=\alpha\), write
\(\overline{\atsign}:=\nu_A\circ\atsign\), and evaluate the resulting
closed expressions using \(\atsign_{\varnothing,\mathcal A^+}\). The same two maps
satisfy the calculation rules

\[
\boxed{
\atsign_{\mathcal A}\,t
=\atsign_{\varnothing,\mathcal A^+}\bigl(t[\overline{\atsign}]\bigr),
\qquad
\varphi\,\atsign_{\mathcal A}
\ \Longleftrightarrow\ 
\varphi[\overline{\atsign}]\,\atsign_{\varnothing,\mathcal A^+}.
}
\]

Thus the bar over \(\atsign\) changes an assignment into a substitution of
names, while the prefix or postfix position of the assignment symbol marks
term evaluation or satisfaction. The right-hand expressions are evaluated in
the distinguished expanded structure \(\mathcal A^+\).

\paragraph{Traditional notation.}
Ordinary evaluation brackets and satisfaction notation express the two rules
as follows:

\[
\llbracket t\rrbracket^{\mathcal A}_\alpha
=\llbracket t[\bar\alpha]\rrbracket^{\mathcal A^+}_{\epsilon_A},
\qquad
\mathcal A,\alpha\models\varphi
\ \Longleftrightarrow\ 
\mathcal A^+,\epsilon_A\models\varphi[\bar\alpha].
\]

Assignment modification is also represented syntactically. For \(x\in X\) and
\(c\in A\),
\(\overline{\alpha[x\mapsto c]}=\bar\alpha[x\mapsto\bar c]\). For the
following equations, write \(t[x:=\bar c]\) for \(j_X(t)[x:=\bar c]\) and
\(\varphi[x:=\bar c]\) for \(\iota_X(\varphi)[x:=\bar c]\), regarding the
ground constant \(\bar c\) as a term in the variable context \(X\). The
structures \(\mathcal A\) and \(\mathcal A^+\) have the same domain, so the
same function \(\alpha:X\to A\) is an assignment in both. The substitution
lemmas yield

\[
\begin{aligned}
\llbracket t\rrbracket^{\mathcal A}_{\alpha[x\mapsto c]}
&=\llbracket t[x:=\bar c]\rrbracket^{\mathcal A^+}_\alpha,\\
\mathcal A,\alpha[x\mapsto c]\models\varphi
&\quad\Longleftrightarrow\quad
\mathcal A^+,\alpha\models\varphi[x:=\bar c].
\end{aligned}
\]

In the \(\atsign\)-calculus, where the structure subscript distinguishes
evaluation in \(\mathcal A\) from evaluation in \(\mathcal A^+\), these are

\[
\boxed{
(\atsign\vert_c^x)_{\mathcal A}\,t
=\atsign_{\mathcal A^+}(\vert_{\bar c}^x t),
\qquad
\varphi\,(\atsign\vert_c^x)_{\mathcal A}
\ \Longleftrightarrow\ 
(\varphi\vert_{\bar c}^x)\,\atsign_{\mathcal A^+}.
}
\]

The expansion therefore internalizes assignments as substitutions of
constant names. It supplies a chosen ground-term name for every element, and
the algebra map evaluates those terms back into the structure.

\paragraph{Example: replacing an assignment by names.}
Let \(\mathbb N^+\) add a constant \(\bar n\) for every natural number.
Take \(X=\{x,y,z\}\), \(t=x\cdot(y+1)\),
\(\varphi=(t=z)\), and an assignment \(\atsign:X\to\mathbb N\) with
\(\atsign(x)=2\), \(\atsign(y)=3\), and \(\atsign(z)=8\). Categorically,
\(\overline{\atsign}=\nu_{\mathbb N}\circ\atsign\) is the ground-term
substitution with
\(t[\overline{\atsign}]=\bar2\cdot(\bar3+1)\) and
\(\varphi[\overline{\atsign}]=
(\bar2\cdot(\bar3+1)=\bar8)\). The corresponding \(\atsign\)-calculation is

\[
\atsign_{\mathbb N}\,t
=\atsign_{\varnothing,\mathbb N^+}\bigl(\bar2\cdot(\bar3+1)\bigr)
=2\cdot(3+1)=8,
\qquad
\varphi\,\atsign_{\mathbb N}
\ \Longleftrightarrow\ 
(\bar2\cdot(\bar3+1)=\bar8)\,\atsign_{\varnothing,\mathbb N^+}.
\]

The ordinary textbook statement is
\(\mathbb N,\atsign\models x\cdot(y+1)=z\) if and only if the sentence
\(\bar2\cdot(\bar3+1)=\bar8\) is true in \(\mathbb N^+\). The categorical
map and the \(\atsign\)-calculation express the same replacement of an
assignment by constant names.

\section{Further perspective: functorial models and generalized points}
\label{functorial-models}

This section may be skipped on a first reading. The functorial reformulation
extends term algebras, assignment actions, and satisfaction from set-based
models to models valued in other categories.

The action of substitutions on assignments already defines a functor. Fix the
functional structure \((A,a)\). Assign to each variable set \(X\) its set
\(A^X\) of assignments, and send a Kleisli substitution
\(\sigma:X\to T_\Sigma Y\) to the map
\(P_a(\sigma):A^Y\to A^X\) defined earlier.

\Needspace{11\baselineskip}
For
\(\sigma:X\to T_\Sigma Y\) and \(\tau:Y\to T_\Sigma Z\), the composition
and identity laws are typed by the following commutative triangles:

\[
\begin{tikzcd}[row sep=large,column sep=huge]
A^Z \arrow[r,"P_a(\tau)"]
  \arrow[dr,"P_a(\tau\star\sigma)"'] &
A^Y \arrow[d,"P_a(\sigma)"] \\
& A^X.
\end{tikzcd}
\qquad
\begin{tikzcd}[row sep=large,column sep=large]
& A^X \arrow[dl,"P_a(\eta_X)"'] \arrow[dr,"1_{A^X}"] & \\
A^X \arrow[rr,"1_{A^X}"'] && A^X.
\end{tikzcd}
\]

These laws make
\(P_a:\mathbf{Kl}(T_\Sigma)^{\mathrm{op}}\to\mathbf{Set}\) a contravariant
functor, with
\(P_a(X)=A^X\). The algebra law therefore organizes assignments over all
variable contexts at once.

For a finitary term monad, the assignment functor is the set-valued instance of
Lawvere's functorial semantics \cite{Lawvere1963}. Let \(\mathbb L_\Sigma\) be the full
subcategory of \(\mathbf{Kl}(T_\Sigma)^{\mathrm{op}}\) on finite cardinals.
Its arrows satisfy
\(\mathbb L_\Sigma(n,m)\cong(T_\Sigma n)^m\): an arrow \(n\to m\) is an
\(m\)-tuple of terms in \(n\) variables. This is the \textbf{Lawvere theory}
of \(\Sigma\). Finite products generalize Cartesian products by their universal
mapping property. A model in a category \(\mathcal C\) with finite products is
a product-preserving functor \(M:\mathbb L_\Sigma\to\mathcal C\); that is, it
sends the abstract product diagrams to product diagrams. Natural
transformations between models are homomorphisms. Taking
\(\mathcal C=\mathbf{Set}\) recovers the \(T_\Sigma\)-algebras used above;
for an equational theory, one instead quotients the arrows by provable equality
of terms.

Changing the target category produces further examples. The Lawvere
theory of groups yields ordinary groups in \(\mathbf{Set}\), topological
groups in \(\mathbf{Top}\), sheaves of groups in \(\mathbf{Sh}(S)\), and
group schemes over \(S\) in \(\mathbf{Sch}/S\).

\Needspace{8\baselineskip}
For a homomorphism
\(h:G\to H\), naturality at multiplication is the familiar commutative
square

\[
\begin{tikzcd}[row sep=large,column sep=huge]
G\times G \arrow[r,"m_G"] \arrow[d,"h\times h"'] &
G \arrow[d,"h"] \\
H\times H \arrow[r,"m_H"'] & H .
\end{tikzcd}
\]

Its commutativity is the equation \(h(g_1g_2)=h(g_1)h(g_2)\). A morphism
of theories induces a reduct functor between model categories by
precomposition, while a finite-product-preserving functor between target
categories transports models by postcomposition.

\subsection{Relations, geometric logic, and classifying toposes}

A Lawvere theory describes operations and equations. Relation symbols and the
operations of geometric logic---finite conjunction, arbitrary disjunction,
and existential quantification---require a richer syntactic category. For a
geometric theory \(\mathbb T\), its classifying topos
\(\mathcal E_{\mathbb T}\) and universal model \(U_{\mathbb T}\) have the
following classifying property~\cite[pp.~3--5]{CaramelloClassifying}:

\[
\operatorname{Geom}(\mathcal E,\mathcal E_{\mathbb T})
\simeq
\mathbb T\text{-}\operatorname{Mod}(\mathcal E),
\qquad f\longmapsto f^*U_{\mathbb T}.
\]

The equivalence is natural in every Grothendieck topos \(\mathcal E\). A
\textbf{Grothendieck topos} is a category of sheaves on a site, up to
equivalence; a site is a category with a notion of covering. A
\textbf{geometric morphism} \(f:\mathcal E\to\mathcal F\) is an adjunction
\(f^*\dashv f_*\) with \(f^*\) preserving finite limits.

\Needspace{11\baselineskip}
For a geometric morphism
\(g:\mathcal E'\to\mathcal E\), naturality means that the following square
commutes up to natural isomorphism:

\[
\begin{tikzcd}[row sep=large,column sep=large]
\operatorname{Geom}(\mathcal E,\mathcal E_{\mathbb T})
  \arrow[r,"f\mapsto f^*U_{\mathbb T}"]
  \arrow[d,"(-)\circ g"'] &
\mathbb T\text{-}\operatorname{Mod}(\mathcal E)
  \arrow[d,"g^*"] \\
\operatorname{Geom}(\mathcal E',\mathcal E_{\mathbb T})
  \arrow[r,"f\mapsto f^*U_{\mathbb T}"'] &
\mathbb T\text{-}\operatorname{Mod}(\mathcal E').
\end{tikzcd}
\]

A \textbf{point} of
a topos \(\mathcal F\) is a geometric morphism
\(\mathbf{Set}\to\mathcal F\). Thus ordinary set-based models correspond,
up to the displayed equivalence, to points of \(\mathcal E_{\mathbb T}\);
models internal to another topos correspond to generalized points. This
statement applies directly to geometric theories. A classical first-order
theory can first be Morleyized to a coherent, hence geometric, theory while
retaining the same set-based models up to canonical expansion and reduct;
internal models then have the semantics of that Morleyization.

The satisfaction pairing also generalizes. Let \(A\) be the carrier of an
internal \(\mathbb T\)-model in a Grothendieck topos \(\mathcal E\), let
\(V\) be an object of \(\mathcal E\), and put \(X=\{1,\ldots,n\}\). A
generalized assignment is an arrow \(\atsign:V\to A^n\). A geometric formula
\(\varphi\) in the context \(X\) is interpreted by a subobject
\(\llbracket\varphi\rrbracket\hookrightarrow A^n\), and its locus of
satisfaction is the pullback \(V_\varphi\hookrightarrow V\):

\[
\begin{tikzcd}[row sep=large,column sep=huge]
V_\varphi \arrow[r] \arrow[d,hook']
  \arrow[dr,phantom,"\lrcorner",very near start] &
\llbracket\varphi\rrbracket \arrow[d,hook] \\
V \arrow[r,"\atsign"'] & A^n .
\end{tikzcd}
\]

For \(\mathcal E=\mathbf{Set}\) and \(V=1\), the subobject
\(V_\varphi\hookrightarrow1\) is either false or true, recovering the
\(\mathbf2\)-valued pairing used earlier. In a general topos, the satisfaction
locus is a subobject of the stage \(V\); global truth values are arrows
\(1\to\Omega_{\mathcal E}\).

\subsection{A spectrum example}

The comparison with algebraic geometry can be made exact. Let \(R\) be a
commutative ring. Introduce one proposition \(D(f)\) for each \(f\in R\) and
impose the geometric sequents \cite{Manuell2022}

\[
D(0)\vdash\bot,\qquad
\top\vdash D(1),\qquad
D(fg)\dashv\vdash D(f)\wedge D(g),\qquad
D(f+g)\vdash D(f)\vee D(g).
\]

\Needspace{9\baselineskip}
A set-based model selects the elements \(f\) for which \(D(f)\) holds. This
set is a prime anti-ideal; classically, its complement is a prime ideal. The
classifying locale of this propositional geometric theory is the Zariski locale
underlying \(\operatorname{Spec}R\); its classifying topos is the associated
localic topos \(\mathbf{Sh}(\operatorname{Spec}R)\). A \textbf{locale} describes a
space through its lattice of open sets. Thus, when
\(R\) is the coordinate
ring of an affine curve, set-based models correspond bijectively, classically,
to the scheme-theoretic points of the curve, including its non-closed points. A prime
\(\mathfrak p\) also determines the field-valued point
\(\operatorname{Spec}\kappa(\mathfrak p)\to\operatorname{Spec}R\)
\cite{StacksPoints}. This is an instance of models as points.

\Needspace{12\baselineskip}
\section{Ultraproducts, stalks, and the
\texorpdfstring{\(\atsign\)}{@}-calculus}
\label{ultraproducts-at-calculus}

The distinction between term evaluation and satisfaction persists under
ultraproducts. Prefix \(\atsign\) calculates a term's value in the ultraproduct;
postfix \(\atsign\) expresses the transfer of satisfaction in
\L{}o\'s's theorem. The interpretation of an ultraproduct as a sheaf stalk
then realizes satisfaction as membership in a clopen subset of an
extremally disconnected space. These constructions also distinguish
evaluation of a fixed finite expression from the formation of an
ultraproduct of syntax sets.

\subsection{Assignments and term evaluation}

Work in classical set theory with choice. Let \(I\) be a set, and let
\(\mathcal U\) be an ultrafilter on \(I\): a collection of subsets closed
under supersets and finite intersections, containing \(I\) and excluding
\(\varnothing\), such that exactly one of \(B\) and \(I\setminus B\) belongs
to \(\mathcal U\) for each \(B\subseteq I\).
Let \((\mathcal A_i)_{i\in I}\) be nonempty structures in one finitary
first-order language \(\mathcal L\), with functional signature \(\Sigma\)
and underlying sets \(A_i\).

Put \(C=\prod_{i\in I}A_i\). For \(a,b\in C\), define
\(a\sim_{\mathcal U}b\) when
\(\{i:a_i=b_i\}\in\mathcal U\). The ultraproduct
\(\mathcal A_{\mathcal U}=\prod_{\mathcal U}\mathcal A_i\) has underlying
set \(A_{\mathcal U}=C/{\sim_{\mathcal U}}\). Write
\(q_{\mathcal U}:C\to A_{\mathcal U}\) for the quotient map and
\([a]_{\mathcal U}=q_{\mathcal U}(a)\). Function symbols act on
representatives coordinatewise; a relation holds of classes precisely when
it holds of their representatives on a member of \(\mathcal U\).
The ultrafilter's closure under finite intersections makes these
interpretations independent of representatives.

Fix a finite set \(X\) of variables and assignments
\(\atsign_i:X\to A_i\). Their product is the assignment
\(\atsign_C:X\to C\) given by
\(\atsign_C(x)=(\atsign_i(x))_{i\in I}\). The induced ultraproduct
assignment is
\(\atsign_{\mathcal U}=q_{\mathcal U}\circ\atsign_C:X\to A_{\mathcal U}\).
For \(t\in T_\Sigma X\), put
\(e(t)=(\atsign_i\,t)_{i\in I}\in C\).

\Needspace{11\baselineskip}
The term-evaluation triangle is
\[
\begin{tikzcd}[row sep=large,column sep=huge]
T_\Sigma X \arrow[r,"e"]
  \arrow[dr,"\llbracket-\rrbracket^{\mathcal A_{\mathcal U}}_{\atsign_{\mathcal U}}"'] &
C \arrow[d,"q_{\mathcal U}"] \\
& A_{\mathcal U}.
\end{tikzcd}
\]
Its calculational form is
\[
\boxed{\atsign_{\mathcal U}\,t
=[(\atsign_i\,t)_{i\in I}]_{\mathcal U}.}
\]
Indeed, the equation for a variable is the definition of
\(\atsign_{\mathcal U}\). If it holds for \(t_1,\ldots,t_n\), the
coordinatewise interpretation of \(f\in\Sigma\) proves it for
\(f(t_1,\ldots,t_n)\). This is structural induction on \(t\), including
the nullary case. Equivalently, \(q_{\mathcal U}\) is a homomorphism of
the functional reducts, and homomorphisms commute with term evaluation.

\subsection{Satisfaction, witnesses, and substitution}

For \(\varphi\in\operatorname{Form}_{\mathcal L}(X)\), define its set of
satisfying indices by
\(B_\varphi=\{i\in I:\varphi\,\atsign_i\}\).
Let \(u_{\mathcal U}:\mathcal P(I)\to\mathbf2\) take the value
\(\top\) exactly on the members of \(\mathcal U\). It preserves finite
Boolean operations: intersections represent conjunction, complements
represent negation, and unions represent disjunction.

\Needspace{12\baselineskip}
\L{}o\'s's theorem~\cite[Corollary~10]{LurieUltraproducts} states that the
following square commutes:
\[
\begin{tikzcd}[row sep=large,column sep=huge]
\operatorname{Form}_{\mathcal L}(X)
  \arrow[r,"\varphi\mapsto B_\varphi"]
  \arrow[d,"S^{\mathcal A_{\mathcal U}}_{\atsign_{\mathcal U}}"'] &
\mathcal P(I) \arrow[d,"u_{\mathcal U}"] \\
\mathbf2 \arrow[r,"1_{\mathbf2}"'] & \mathbf2.
\end{tikzcd}
\]
In the postfix notation, this is
\[
\boxed{\varphi\,\atsign_{\mathcal U}
\quad\Longleftrightarrow\quad
\{i\in I:\varphi\,\atsign_i\}\in\mathcal U.}
\]

Here is a proof in the assignment calculus. Atomic equality follows from
the term-evaluation identity and the definition of
\(\sim_{\mathcal U}\); the other atomic cases follow from the
interpretation of relation symbols. Finite Boolean connectives follow
from the corresponding properties of \(u_{\mathcal U}\).

For the existential step, suppose
\(B=\{i:((\exists x)\psi)\,\atsign_i\}\in\mathcal U\).
Choose \(b_i\in A_i\) such that
\(\psi(\atsign_i\vert_{b_i}^{x})\) for each \(i\in B\), and choose an
arbitrary \(b_i\in A_i\) elsewhere. Put \(b=[(b_i)]_{\mathcal U}\).
The updated assignments \(\atsign_i\vert_{b_i}^{x}\), on the context
\(X\cup\{x\}\), induce exactly
\(\atsign_{\mathcal U}\vert_b^{x}\). The induction hypothesis applied to
\(\psi\) therefore yields
\(\psi(\atsign_{\mathcal U}\vert_b^{x})\), and hence
\(((\exists x)\psi)\,\atsign_{\mathcal U}\).
Conversely, represent a witness \(b\in A_{\mathcal U}\) by a tuple
\((b_i)\). The induction hypothesis says that
\(\{i:\psi(\atsign_i\vert_{b_i}^{x})\}\in\mathcal U\); this set is
contained in \(B\), so \(B\in\mathcal U\).
Universal quantification follows from
\((\forall x)\psi\equiv\neg(\exists x)\neg\psi\).

The same update calculation combines transfer with the substitution
lemma. For \(x\in X\), \(s\in T_\Sigma X\), and capture-avoiding
substitution, it yields
\[
\begin{aligned}
(\varphi\vert_s^x)\,\atsign_{\mathcal U}
&\quad\Longleftrightarrow\quad
\varphi\bigl(\atsign_{\mathcal U}
                  \vert_{\atsign_{\mathcal U}\,s}^{x}\bigr)\\
&\quad\Longleftrightarrow\quad
\bigl\{i:\varphi\bigl(\atsign_i\vert_{\atsign_i\,s}^{x}\bigr)\bigr\}
\in\mathcal U.
\end{aligned}
\]
The first step moves syntactic substitution to the assignment; the second
applies \L{}o\'s's theorem to the family of updated assignments.

\Needspace{12\baselineskip}
\subsection{Ultraproducts as stalks and truth as a locus}

Regard \(I\) as a discrete space. Its Stone--\v{C}ech compactification
\(\beta I\) is the space of ultrafilters on \(I\), with clopen basis
\(\widehat B=\{\mathcal V:B\in\mathcal V\}\), for \(B\subseteq I\).
The map \(B\mapsto\widehat B\) identifies \(\mathcal P(I)\) with the
Boolean algebra of clopen subsets of \(\beta I\).
An index \(i\) corresponds to the principal ultrafilter
\(\{B\subseteq I:i\in B\}\).

Let \(j:I\hookrightarrow\beta I\) be this inclusion.
The family \((A_i)\) is a sheaf on \(I\), whose sections over
\(B\subseteq I\) are \(\prod_{i\in B}A_i\).
Its direct image \(\mathcal F\) on \(\beta I\) satisfies
\(\mathcal F(\widehat B)=\prod_{i\in B}A_i\), with restriction maps
given by coordinate projections. Consequently,
\[
\mathcal F_{\mathcal U}
=\varinjlim_{B\in\mathcal U}\prod_{i\in B}A_i
\cong A_{\mathcal U},
\]
where the colimit is taken along restriction to smaller members of
\(\mathcal U\). Two local sections have the same germ precisely when
they agree after restriction to a common member of \(\mathcal U\).
Nonemptiness and choice allow a tuple on such a member to extend to
all of \(I\), giving the displayed identification with the quotient
ultraproduct~\cite[Example~14]{LurieUltraproducts}.

The assignments \(\atsign_i\) determine
\(\atsign_\beta:1\to\mathcal F^X\) in
\(\mathbf{Sh}(\beta I)\), where \(1\) is the terminal sheaf.
Its stalk at \(\mathcal U\) is \(\atsign_{\mathcal U}\).
Term evaluation yields the global section
\(\atsign_\beta\,t=(\atsign_i\,t)_i\).
The germ map \(g_{\mathcal U}:\mathcal F(\beta I)\to\mathcal F_{\mathcal U}\)
therefore satisfies
\(g_{\mathcal U}(\atsign_\beta\,t)=\atsign_{\mathcal U}\,t\).
Prefix \(\atsign\) returns a section before taking the stalk and an
element afterwards.

For the chosen assignments, \L{}o\'s's theorem also identifies the locus
of ultrafilters satisfying a fixed first-order formula:
\[
\{\mathcal V\in\beta I:\varphi\,\atsign_{\mathcal V}\}
=\widehat{B_\varphi}.
\]
This has the internal satisfaction-locus interpretation for coherent
formulas, built from atomic formulas using finite conjunctions, finite
disjunctions, and existential quantification. Interpret the structures
as a sheaf model \(\mathcal F\); let
\(\llbracket\varphi\rrbracket_{\mathcal F}\hookrightarrow\mathcal F^X\)
be the interpreted relation. For an open subset \(O\subseteq\beta I\),
write \(1_O\hookrightarrow1\) for its associated subterminal sheaf:
its sections over an open \(W\) form a singleton if \(W\subseteq O\)
and are empty otherwise.
Then the following square is a pullback:

\Needspace{10\baselineskip}
\[
\begin{tikzcd}[row sep=large,column sep=huge]
1_{\widehat{B_\varphi}} \arrow[r] \arrow[d,hook']
  \arrow[dr,phantom,"\lrcorner",very near start] &
\llbracket\varphi\rrbracket_{\mathcal F} \arrow[d,hook] \\
1 \arrow[r,"\atsign_\beta"'] & \mathcal F^X.
\end{tikzcd}
\]
Indeed, direct image along \(j\) preserves finite limits, finite
coproducts, and images in this situation, so it preserves coherent
interpretation~\cite[Corollary~1.10]{HaineCondensedPoints}.
For arbitrary first-order formulas, the same description applies to
their predicates in the canonical coherent Morleyization.
At a general stage \(V\) of an internal model, satisfaction remains a
subobject \(V_\varphi\hookrightarrow V\); at a set-valued point it
reduces to the postfix assertion \(\varphi\,\atsign\).

The space \(\beta I\) is compact Hausdorff and extremally disconnected:
closures of open subsets are open.
Such spaces form a defining site for condensed mathematics. With size
bounds fixed, condensed sets can be described as contravariant functors
on these spaces that take finite disjoint unions to
products~\cite[Example~2.5 and Proposition~2.8]{ScholzeCondensed}.
For a coherent theory \(\mathbb T\), the analogous assignment
\[
\mathbf{Mod}_{\mathbb T}(K)
=\mathbb T\text{-}\operatorname{Mod}\bigl(\mathbf{Sh}(K)\bigr)
\]
defines a condensed category of models: inverse image determines its
contravariant dependence on \(K\), and disjoint unions of spaces yield
products of model categories. Its value at a point is the ordinary
category of set-based models. An indexed family of \(\mathbb T\)-models
extends to an object
over \(\beta I\), and its value at an ultrafilter point is the
corresponding ultraproduct~\cite[Example~0.2 and Proposition~1.6]{HaineCondensedPoints}.
Here ordinary sheaves on each \(K\) supply the values of the condensed
category as \(K\) varies.

\subsection{Finite fields and an infinite disjunction}

Let \(P\) be the set of primes and take \(\mathcal A_p=\mathbb F_p\).
Use \(\atsign_p:\varnothing\to\mathbb F_p\) for the empty assignment.
For \(m\geq1\), let
\(\chi_m=(m\cdot1=0)\), where \(m\cdot1\) abbreviates the term formed
by adding \(1\) to itself \(m\) times. Then
\(B_{\chi_m}=\{p\in P:p\mid m\}\) is finite.

A nonprincipal ultrafilter contains no finite set: if a finite union of
singletons belongs to an ultrafilter, one of those singletons belongs
to it, making it principal. Hence
\(\neg(\chi_m\,\atsign_{\mathcal U})\) for every \(m\geq1\) and every
nonprincipal \(\mathcal U\) on \(P\).
The ultraproduct is a field by \L{}o\'s's theorem, and these inequalities
show that its characteristic is zero.

Now form the infinitary geometric sentence
\(\psi=\bigvee_{m\geq1}\chi_m\).
Every \(\psi\,\atsign_p\) is true, since \(\chi_p\,\atsign_p\) is true.
But \(\psi\,\atsign_{\mathcal U}\) is false at every nonprincipal
ultrafilter: none of its disjuncts holds there.
The first-order transfer theorem does not include this infinite
disjunction.

The sheaf interpretation makes the difference explicit.
Each \(\widehat{B_{\chi_m}}\) consists of the finitely many principal
points indexed by primes dividing \(m\). The satisfaction locus of the
geometric disjunction in \(\mathbf{Sh}(\beta P)\) is therefore
\[
\bigcup_{m\geq1}\widehat{B_{\chi_m}}=P\subset\beta P.
\]
This open subset is dense and proper, and so is not clopen. Density
follows because every nonempty basic clopen meets \(P\); properness
follows from the existence of nonprincipal ultrafilters.
Taking a stalk preserves the geometric disjunction, which is false at
every point of \(\beta P\setminus P\).
The direct-image functor \(j_*\), however, does not preserve this
infinite disjunction: the disjunction is everywhere true on \(P\),
whose terminal sheaf has direct image the terminal sheaf on \(\beta P\).
Thus a satisfaction locus can carry more information than a single
external truth value.

\subsection{Fixed terms and ultraproducts of syntax}

The prefix evaluation identity uses one fixed term
\(t\in T_\Sigma X\). The comparison between forming terms and forming
ultraproducts has a different type:
\[
\lambda_{\mathcal U}:
T_\Sigma\!\left(\prod_{\mathcal U}A_i\right)
\longrightarrow
\prod_{\mathcal U}T_\Sigma A_i.
\]
To construct it, give the target its coordinatewise
\(\Sigma\)-algebra structure and send a variable
\([a]_{\mathcal U}\) to
\([(\eta_{A_i}(a_i))_i]_{\mathcal U}\).
This variable map is well-defined, since representatives agreeing on
a member of \(\mathcal U\) have corresponding variable terms agreeing
there. The free-algebra property extends it uniquely to
\(\lambda_{\mathcal U}\).
Since a term has only finitely many variable occurrences, its image can
be computed by choosing representatives of their labels in
\(A_{\mathcal U}\) and keeping the same term shape in every coordinate.

The comparison need not be surjective. Let \(I=\mathbb N\), let
\(\mathcal U\) be nonprincipal, let every \(A_i=\{*\}\), and let
\(\Sigma\) consist of one unary symbol \(f\). Then
\(T_\Sigma(\prod_{\mathcal U}A_i)=T_\Sigma\{*\}\) consists of the terms
\(f^n(*)\), for \(n\in\mathbb N\), and
\(\lambda_{\mathcal U}(f^n(*))=[(f^n(*))_i]_{\mathcal U}\).
The target also contains
\[
\tau=[(f^i(*))_{i\in\mathbb N}]_{\mathcal U}.
\]
For each fixed \(n\), the equality \(f^i(*)=f^n(*)\) holds only when
\(i=n\), so \(\tau\neq\lambda_{\mathcal U}(f^n(*))\).
Its representatives have unbounded term depth.
Thus the ultraproduct of term sets contains elements that are not
represented by any ordinary finite term over the ultraproduct.
This does not affect the evaluation identity for a fixed term.

\subsection{Names in a common expanded language}

The expansion by names also applies, using a common signature for all
factors. Retain \(C=\prod_i A_i\), and add a constant \(\bar c\) for
each \(c\in C\), obtaining \(\mathcal L_C\).
Interpret \(\bar c\) in \(\mathcal A_i^C\) as \(c_i\), and in
\(\mathcal A_{\mathcal U}^C\) as \([c]_{\mathcal U}\).
These are the expansions determined by the maps
\(C\to A_i\) and \(q_{\mathcal U}:C\to A_{\mathcal U}\), respectively.
The expanded ultraproduct is precisely
\(\mathcal A_{\mathcal U}^C=\prod_{\mathcal U}\mathcal A_i^C\).

Write \(\atsign_{\varnothing,\mathcal A_i^C}:\varnothing\to A_i\) and
\(\atsign_{\varnothing,\mathcal A_{\mathcal U}^C}:\varnothing\to A_{\mathcal U}\)
for the empty assignments in these expanded structures. The name-evaluation calculation is
\[
\atsign_{\varnothing,\mathcal A_{\mathcal U}^C}\,\bar c
=[(\atsign_{\varnothing,\mathcal A_i^C}\,\bar c)_i]_{\mathcal U}
=[c]_{\mathcal U}.
\]
Different names \(\bar c,\bar d\) have the same value precisely when
\(c\sim_{\mathcal U}d\); they remain different symbols when \(c\neq d\).

The product assignment \(\atsign_C:X\to C\) determines the common
ground-term substitution
\(\overline{\atsign_C}=\nu_C\circ\atsign_C:
X\to T_{\Sigma[C]}\varnothing\).
Its value at \(x\) is the name of the tuple
\((\atsign_i(x))_i\).
For an original term \(t\) and formula \(\varphi\), regarded in
\(\mathcal L_C\), the naming identities are
\[
\begin{aligned}
\atsign_{\mathcal U}\,t
&=\atsign_{\varnothing,\mathcal A_{\mathcal U}^C}
       \bigl(t[\overline{\atsign_C}]\bigr),\\
\varphi\,\atsign_{\mathcal U}
&\quad\Longleftrightarrow\quad
\varphi[\overline{\atsign_C}]\,
       \atsign_{\varnothing,\mathcal A_{\mathcal U}^C}.
\end{aligned}
\]
Evaluating the inserted name for \(x\) yields
\([(\atsign_i(x))_i]_{\mathcal U}=\atsign_{\mathcal U}(x)\);
the two identities therefore follow from the term and formula
substitution lemmas.

\subsection{Occurrence maps and transferred truth labels}

Fix the finite occurrence-substitution map
\(\gamma:=\gamma_{M,\vec o}:D_{M,\vec n}\to C_{\vec\ell}\) of
Section~\ref{occurrence-indexed-substitution}. Let
\(e_i:C_{\vec\ell}\to\mathbf2\) label the formal arguments by
their truth values in \(\mathcal A_i\) under \(\atsign_i\), and let
\(d_i=e_i\circ\gamma\) label their occurrences.
Define \(e_{\mathcal U}\) and \(d_{\mathcal U}\) by
\[
\begin{aligned}
e_{\mathcal U}(c)=\top
&\quad\Longleftrightarrow\quad
\{i:e_i(c)=\top\}\in\mathcal U,\\
d_{\mathcal U}(r)=\top
&\quad\Longleftrightarrow\quad
\{i:d_i(r)=\top\}\in\mathcal U.
\end{aligned}
\]
For every \(r\), the two sets
\(\{i:d_i(r)=\top\}\) and
\(\{i:e_i(\gamma(r))=\top\}\) are equal. Hence
\(d_{\mathcal U}=e_{\mathcal U}\circ\gamma\).
When \(e_i(c)=\top\) means
\(\psi_c\,\atsign_i\) for a fixed formula \(\psi_c\),
\L{}o\'s's theorem identifies
\(e_{\mathcal U}(c)=\top\) with
\(\psi_c\,\atsign_{\mathcal U}\).
Thus the occurrence map transports exactly the truth labels supplied
by postfix \(\atsign\) in the ultraproduct. Evaluation of the fixed
finite connective pattern then commutes with
\(u_{\mathcal U}\).


\Needspace{12\baselineskip}
\section{Conclusion}\label{conclusion}

The principal objects and maps that elementary notation can obscure are now
separated by type:

\[
\begin{array}{rcl}
\text{terms over }X &:& T_\Sigma X,\\
\text{syntactic substitutions} &:& X\to T_\Sigma Y,\\
\text{assignments} &:& X\to A,\\
\text{term evaluation} &:& A^X\times T_\Sigma X\to A,\\
\text{formulas over }X &:& \operatorname{Form}_{\mathcal L}(X),\\
\text{satisfaction} &:&
\operatorname{Form}_{\mathcal L}(X)\times A^X\to\mathbf2.
\end{array}
\]

Categorically, substitution is composition in \(\mathbf{Kl}(T_\Sigma)\),
evaluation is the action of an Eilenberg--Moore algebra, and the substitution
lemma is their compatibility law. Language expansion turns assignments into
ground-term substitutions, with the unit law ensuring that each name evaluates
to the element it names. The \DP\ occurrence maps record generalized
composition of formula patterns by tracking repeated argument occurrences
through finite coproduct fibres and their ordinal encodings.

The parallel $\atsign$-calculus writes the same maps as calculations:
$\atsign\,t$ returns a term's value, \(\varphi\,\atsign\) records satisfaction,
and the sliding identities calculate compatibility with substitution. The two
presentations have complementary jobs: categorical maps expose types and compositional
laws, while the $\atsign$-calculus supplies the working notation.

The categorical organization is also functorial. Lawvere theories admit models
in any category with finite products, and classifying toposes represent models
of geometric theories by points or generalized points. Thus the set-based
account is one instance of the construction.

\clearpage
\appendix
\numberwithin{equation}{section}
\section{Feferman's arithmetization conventions}\label{app:feferman}

Arithmetization adds another distinction to those of the preceding sections:
an expression, its numerical code, and an arithmetic term denoting that code
have different types. Numeral substitution combines all three. The
evaluation identity in Appendix~\ref{app:numeral-substitution} connects this
construction with the prefix $\atsign$-calculus.

In the formula displays of this appendix, $\varphi(v)$ designates a formula
with the indicated free-variable slot, following the textbook convention
described in Section~\ref{formulas-and-satisfaction}; $\varphi(\bn n)$
denotes its numeral instance. Dotted occurrences inside quotation marks
will be defined by an explicit translation into arithmetic terms.

Feferman represents operations on formal syntax by terms and formulas
of arithmetic. We explain the conventions in his \emph{Arithmetization
of metamathematics in a general setting} (1960), especially \S4,
pp.~57--58 \cite{F60}, and \emph{Transfinite recursive progressions of
axiomatic theories} (1962), \S2(c), pp.~268--270 \cite{F62}.

His notation identifies expressions with their G\"odel numbers and
marks represented operations by dots. A dotted variable abbreviates
the represented numeral constructor; combining it with substitution
produces the code of a numeral instance. The exposition distinguishes
expressions, their numerical codes, and the arithmetic terms denoting
those codes. The explicit grammar and coding examples below support
this reading of Feferman's notation. The final section explains its
use in Di Paola and Montagna's parameterized consistency formulas.

\subsection{Numbers, numerals, and the function \texorpdfstring{$\Num$}{Num}}\label{app:numerals}

\subsubsection*{The arithmetic language}
The first-order language of arithmetic used here has nonlogical signature
\[
 \LA=\{0,S,+,\cdot\},
\]
where $0$ is a constant symbol, $S$ is unary, and $+$ and $\cdot$
are binary function symbols. Equality is a logical symbol. With variables
$v_0,v_1,\ldots$, its terms $t$ and formulas $A$ are generated by
\begin{align*}
 t&::=v_i\mid0\mid S(t)\mid t+t\mid t\cdot t,\\
 A&::=t=t\mid\neg A\mid A\to A\mid\forall v_i\,A.
\end{align*}
Parentheses specify the grouping of expressions. Other connectives and
quantifiers are abbreviations in this syntax. The symbol $\cdot$ is
multiplication; the dot above a variable has a different role.

The corresponding standard structure is
\[
 \mathcal N=\langle\N,0,S^{\mathcal N},+^{\mathcal N},
                         \cdot^{\mathcal N}\rangle,
 \qquad S^{\mathcal N}(n)=n+1,
\]
with addition and multiplication of natural numbers. Thus the signature
lists symbols and their arities, whereas the structure also specifies a
domain and interpretations. Below, $\N$ also denotes this structure and
its expansion by the specified primitive recursive functions.

Fix a language $L^{*}$ extending $\LA$ by function symbols for the
primitive recursive functions needed below. Fix an
injective G\"odel coding $E\mapsto\code{E}$ of its finite expressions,
with primitive recursive syntax predicates and constructors. Expressions
are coded as raw syntax: distinct bound-variable names remain distinct,
although Section~\ref{formulas-and-satisfaction} takes formulas modulo
$\alpha$-equivalence. Provably equal terms and logically equivalent
formulas also retain their distinct codes. All displayed formula abbreviations are
expanded by fixed conventions before coding.

\subsubsection*{Numerals and their codes}
For $n\in\N$, its \emph{numeral} is the closed term defined by
\[
  \bn0\defeq 0,\qquad \bn{n+1}\defeq S(\bn n).
\]
Thus $\bn2$ is the term $S(S(0))$. Throughout this appendix, $\bn n$
denotes this numeral, whereas the bar in
Section~\ref{expansion-of-a-language-by-names-for-a-structure} denotes an
adjoined constant symbol naming an element. No new constant symbols
are needed for these numerals. Define the numerical function
\begin{equation}\label{eq:num}
  \boxed{\num(n)\defeq\code{\bn n}.}
\end{equation}
Its input is a number; its output is another number, namely a code.
If $\scod(\code t)=\code{S(t)}$ for each term $t$, with
$\scod(a)=\code0$ on non-term codes, then
\begin{equation}\label{eq:num-recursion}
  \num(0)=\code0,\qquad
  \num(n+1)=\scod(\num(n)).
\end{equation}
These equations prove that $\num$ is primitive recursive. In particular,
$\num(n+1)$ is obtained by coding an additional successor symbol; it is
not defined as $\num(n)+1$.

Write $\PA$ for a conservative definitional extension of PA containing
the chosen primitive recursive function symbols and their defining
equations. The symbol $\Num$ represents $\num$, and $\Sc$ represents
$\scod$. Their equations include
\[
  \PA\vdash \Num(0)=\bn{\code0},\qquad
  \PA\vdash \forall x\,\Num(Sx)=\Sc(\Num(x)).
\]
For every standard $n$,
\begin{equation}\label{eq:num-evaluation}
  \PA\vdash\Num(\bn n)=\bn{\num(n)}.
\end{equation}
Lower-case $\num$ denotes the numerical function in the metatheory;
upper-case $\Num$ is the arithmetic function symbol. When this
distinction is suppressed, both are often written $\Num$.

The construction factors through the set of numeral expressions:
\[
\begin{tikzcd}[column sep=large,row sep=large]
 \N \arrow[r,"n\mapsto\bn n"] \arrow[dr,"\num"'] &
 \mathsf{Numerals}(L^{*}) \arrow[d,"\#"]\\
 & \N.
\end{tikzcd}
\]

\noindent The following expressions have different syntactic roles.
\begin{center}
\renewcommand{\arraystretch}{1.18}
\begin{tabularx}{\linewidth}{@{}lX@{}}
\toprule
Expression & Meaning\\
\midrule
$n$ & A natural number in the metatheory.\\
$\bn n$ & An arithmetic term denoting $n$.\\
$\num(n)$ & The natural number $\code{\bn n}$.\\
$\Num(x)$ & An arithmetic term with free variable $x$.\\
$\Num(\bn n)$ & A closed arithmetic term denoting $\num(n)$.\\
$\bn{\num(n)}$ & The numeral denoting $\num(n)$.\\
\bottomrule
\end{tabularx}
\end{center}
For example, $\Num(\bn2)=\Num(S(S(0)))$ and
$\bn{\num(2)}$ have the same value, but they are different terms and
therefore have different G\"odel codes.

Feferman's name for the numerical constructor is $\mathrm{nm}$.
The dot convention for represented functions appears in
\cite[\S4, p.~57]{F60}, with the numeral constructor on p.~58.
The primitive recursive extensions are introduced in
\cite[\S3, Definition~3.3 and Theorem~3.4, pp.~52--53]{F60}.
In \cite[\S2(c), p.~269]{F62}, Feferman explicitly abbreviates the
represented numeral constructor applied to $u$ by $\dot u$.
In our notation,
\begin{equation}\label{eq:dot}
  \dot u\quad\text{abbreviates}\quad\Num(u).
\end{equation}
This is a special case of the general dot convention, not a separate
operation on variables. Feferman's description of the output as the
$(u+1)$-st numeral counts $\bn0$ as the first numeral.

\subsubsection*{Concrete values of \texorpdfstring{$\#$}{the coding map}}
The number $\code E$ depends on the chosen coding. For explicit numerical
examples, use the following coding of syntax trees. These numerical
assignments are chosen here; they are not Feferman's particular numbering.
The remaining definitions and results apply to any coding with the
primitive recursive operations already specified.

Define the pairing function and the constructor-code function by
\[
 \pi(a,b)\defeq\frac{(a+b)(a+b+1)}{2}+b,
 \qquad c(j,a)\defeq1+\pi(j,a).
\]
The first argument $j$ of $c$ records which syntactic constructor is used.
Code terms and formulas recursively as follows:
\begin{align*}
 \code0&=c(0,0), & \code{v_i}&=c(1,i),\\
 \code{S(t)}&=\ctag2{\code t},
   &\code{t+u}&=\ctag3{\pair{\code t}{\code u}},\\
 \code{t\cdot u}&=\ctag4{\pair{\code t}{\code u}},
   &\code{t=u}&=\ctag5{\pair{\code t}{\code u}},\\
 \code{\neg A}&=\ctag6{\code A},
   &\code{A\to B}&=\ctag7{\pair{\code A}{\code B}},\\
 \code{\forall v_i\,A}&=\ctag8{\pair i{\code A}}.
\end{align*}
Injectivity follows by recovering the constructor and its arguments from
the pairing. Parentheses determine the syntax tree and are not separately
coded. The coding extends to each added function symbol $f_j$ of arity $r$
by the rule
\[
 \code{f_j(t_1,\ldots,t_r)}
   =c\bigl(9+j,\sigma(\code{t_1},\ldots,\code{t_r})\bigr),
\]
where the finite-sequence code satisfies
$\sigma()=0$ and
$\sigma(a_1,\ldots,a_r)=1+\pi(a_1,\sigma(a_2,\ldots,a_r))$ for $r>0$.
The added symbols are indexed with primitive recursive arity lookup.

For this coding the following are actual natural-number values of $\#$:
\begin{center}
\renewcommand{\arraystretch}{1.12}
\begin{tabular}{@{}lr@{}}
\toprule
Expression $E$ & $\code E$\\
\midrule
$0$ & $1$\\
$v_0$ & $2$\\
$S(0)$ & $8$\\
$S(S(0))$ & $64$\\
$0+0$ & $33$\\
$0=0$ & $50$\\
\bottomrule
\end{tabular}
\end{center}
For example,
\[
 \code{S(0)}=1+\pi(2,1)=8,\qquad
 \code{0=0}=1+\pi\bigl(5,\pi(1,1)\bigr)=50.
\]
Consequently $\num(0)=1$, $\num(1)=8$, and $\num(2)=64$;
in particular, $\PA\vdash\Num(\bn2)=\bn{64}$.
The terms $0$ and $0+0$ both denote zero, but their codes are $1$ and $33$.

\subsection{Quotation and substitution}

For a fixed expression $E$, its \emph{ordinary G\"odel quotation} is
the numeral
\begin{equation}\label{eq:quote}
  \quo E\defeq\bn{\code E}.
\end{equation}
It is closed even when $E$ is open. Thus $x$ is not free in
$\quo{x=x}$.
For the concrete coding above,
\[
 \quo{S(0)}=\bn8,\qquad
 \quo{0=0}=\bn{50},\qquad
 \quo{v_0=0}=\bn{86}.
\]
The number $86$ codes the open formula $v_0=0$; its quotation is the
closed numeral $\bn{86}$.

For each $k\geq1$, let
\[
 \sbop_k(a;v_1,\ldots,v_k;b_1,\ldots,b_k)
\]
be the primitive recursive function that simultaneously substitutes
the terms with codes $b_i$ for the free occurrences of the distinct
variables with codes $v_i$ in the expression with code $a$.
Fix a deterministic procedure for renaming bound variables when needed
to avoid capture; preserve the existing names when no renaming is
needed. Set the result to $0$ on invalid inputs. Let $\Sb_k$ be its
representing function symbol in $L^{*}$, and write $\Sb$ for $\Sb_1$.
The semicolons separate blocks of arguments; they are not additional
object-language operations. Feferman's corresponding substitution
operations are $\mathrm{Sb}$ and simultaneous substitution
$\mathrm{SSb}$ \cite[\S2(b), pp.~42--43; \S4, p.~57]{F60}.

For a fixed expression $E$ and a designated variable $v$, define
\begin{equation}\label{eq:q}
 q_{E;v}(t)\defeq
 \Sb\bigl(\quo E;\quo v;\Num(t)\bigr).
\end{equation}
This is an arithmetic term. Its value when $t$ has value $n$ is
\[
 \code{E[\bn n/v]},
\]
where the bracket denotes replacement of the free occurrences of $v$.
The inserted expression is a numeral, so it has no variables that could
be captured. When $E=\varphi(v)$ and $v$ is understood, write
$q_\varphi(t)$ for $q_{\varphi;v}(t)$.

The conventional notation
\begin{equation}\label{eq:dotted-quotation}
 \quo{\varphi(\dot x)}\defeq q_\varphi(x)
   =\Sb\bigl(\quo{\varphi(v)};\quo v;\Num(x)\bigr)
\end{equation}
therefore denotes the code of an instance, as a function of $x$.
In contrast, $\quo{\varphi(v)}$ is the fixed numeral for the code of
the open formula. Feferman's substitution abbreviation on p.~269 of
\cite{F62} is precisely this numeral-substitution construction,
in his convention identifying expressions with their codes.

\subsection{Making numeral substitution explicit}

To make the substitution conventions explicit, the following auxiliary
grammar designates the occurrences to be filled with numerals. Its
translation uses Feferman's represented numeral constructor and
substitution operations. The occurrence syntax is supplied for this
exposition: it distinguishes ordinary variables in the quoted expression
from the arithmetic parameters that determine the inserted numerals.

\subsubsection*{Templates and arithmetic expressions}
Let $v$ range over object variables, and let
$\h_1,\h_2,\ldots$ be distinct \emph{hole symbols}, which are not
variables and cannot be bound. A term template $U$ and formula template
$A$ have grammar
\begin{align*}
 U&::=v\mid0\mid S(U)\mid U+U\mid U\cdot U
        \mid f(U_1,\ldots,U_r)\mid\h_i,\\
 A&::=U=U\mid\neg A\mid A\to A\mid\forall v\,A.
\end{align*}
Here $f$ ranges over the chosen function symbols of $L^{*}$, with
their assigned arities. A template $C$ is either kind. It contains no
quotation constructor. Connectives such as $\land,\lor,\exists$, and
formula notations such as $\Prov_T$ and $\Con_\alpha$, are fixed
abbreviations, expanded before a template is coded.

Surface arithmetic terms $t$ and formulas $B$ have grammar
\begin{align*}
 t&::=x\mid0\mid S(t)\mid t+t\mid t\cdot t
       \mid f(t_1,\ldots,t_r)\mid\Q(C;t_1,\ldots,t_k),\\
 B&::=t=t\mid\neg B\mid B\to B\mid\forall x\,B.
\end{align*}
The distinct holes occurring in $C$ must be exactly
$\h_1,\ldots,\h_k$; a hole may occur more than once. When $k=0$,
write $\Q(C;)$.

\subsubsection*{Translation into the arithmetic language}
Write $t^{\flat}$ for the translation of a surface term into $L^{*}$.
Translate variables and arithmetic constructors component by component.
For $\Q$, choose distinct variables $w_1,\ldots,w_k$ that occur nowhere
in $C$, by a fixed rule. Let $C^{\mathbf w}$ replace each hole
$\h_i$ by $w_i$. Set
\begin{align}
 \flatof{\Q(C;)}&\defeq\quo C,\label{eq:grammar-zero}\\
 \flatof{\Q(C;t_1,\ldots,t_k)}
 &\defeq\Sb_k\bigl(\quo{C^{\mathbf w}};
       \quo{w_1},\ldots,\quo{w_k};
       \Num(t_1^{\flat}),\ldots,\Num(t_k^{\flat})\bigr)
       \quad(k>0).\label{eq:grammar}
\end{align}
Formula translation is likewise componentwise. Equations
\eqref{eq:grammar-zero}--\eqref{eq:grammar} eliminate quotation and holes;
the result is an ordinary arithmetic expression.

The free-variable rule is
\begin{equation}\label{eq:fv}
 \FV\bigl(\flatof{\Q(C;t_1,\ldots,t_k)}\bigr)
   =\bigcup_{i=1}^{k}\FV(t_i^{\flat}).
\end{equation}
Ordinary variables in $C$ occur inside fixed code numerals and contribute
no free variables to the translated term. A quantifier outside $\Q$
can bind variables in its arguments. A quantifier inside $C$ binds
only the ordinary object-variable occurrences in its own scope.

Dotted display notation abbreviates this constructor: put $\dot t_i$
at the occurrences of $\h_i$, and enclose the resulting display in
G\"odel corners. The dot records \emph{numeral insertion} at those
occurrences. Thus a generalized dotted term uses $\Num(t_i^{\flat})$,
not the code of the term $t_i$ itself.

\begin{example}[Ordinary occurrences remain quoted]\label{ex:mixed}
The display $\quo{x=\dot x}$ means $\Q(x=\h_1;x)$. With a fresh $w$,
its translation is
\[
  \Sb(\quo{x=w};\quo w;\Num(x)).
\]
At $x=n$, its value is $\code{x=\bn n}$, an open-formula code.
Replacing the dot by the ordinary variable $x$ and then substituting
for every free $x$ in $x=x$ would instead produce
$\code{\bn n=\bn n}$. Fresh placeholders prevent this error.
\end{example}

\begin{example}[A hole under a quantifier]\label{ex:bound}
The display $\quo{\forall x\,(x=\dot x)}$ means
$\Q(\forall x\,(x=\h_1);x)$. It translates as
\[
  \Sb(\quo{\forall x\,(x=w)};\quo w;\Num(x)),
\]
and has value $\code{\forall x\,(x=\bn n)}$ at $x=n$.
The inner $\forall x$ does not bind the parameter supplied to the
hole. Consequently this grammar needs no prohibition on a dotted $x$
under an inner $\forall x$. The explicit $\Q$ form resolves the two
roles of the same printed letter.
\end{example}

\begin{example}[Repeated holes and compound parameters]
The term $\Q(\h_1=\h_1;x)$ has value
$\code{\bn n=\bn n}$ at $x=n$. The term
$\Q(\h_1=\h_2;x,Sx)$ has value
$\code{\bn n=\bn{n+1}}$. For a fixed $\varphi(v)$,
$\quo{\varphi(\dot{(x+y)})}$ denotes $q_\varphi(x+y)$: at
$x=n,y=m$ it codes $\varphi(\bn{n+m})$, not
$\varphi(\bn n+\bn m)$.
\end{example}

\subsection{The substitution lemma}\label{app:numeral-substitution}

\begin{proposition}[Numeral evaluation]\label{prop:evaluation}
Let the holes of $C$ be $\h_1,\ldots,\h_k$.
For each tuple of standard natural numbers $n_1,\ldots,n_k$,
\begin{equation}\label{eq:lemma}
 \PA\vdash
 \flatof{\Q(C;\bn{n_1},\ldots,\bn{n_k})}
   =\quo{C[\bn{n_1}/\h_1,\ldots,\bn{n_k}/\h_k]}.
\end{equation}
The right side is an ordinary quotation, hence a numeral.
\end{proposition}

\begin{proof}
For $k=0$, the assertion is the defining identity
\eqref{eq:grammar-zero}. For $k>0$, replace each $\Num(\bn{n_i})$ by $\bn{\num(n_i)}$ using
\eqref{eq:num-evaluation}. The remaining closed $\Sb_k$ term
computes the code obtained by substituting $\bn{n_i}$ for the fresh
$w_i$. Those variables occur exactly where the holes occurred and
are not bound. The result is exactly the expression on the right.
The defining equations prove this finite computation in $\PA$.
\end{proof}

In particular, when $q(x)=\flatof{\Q(C;x)}$,
\[
  \PA\vdash q(x)[\bn n/x]=\quo{C[\bn n/\h_1]}.
\]
The substitution on the left is substitution in the translated
arithmetic term. Undotted variable occurrences inside $C$ remain quoted.
More generally, if a valuation $s$ assigns $t_i^{\flat}$ the value $n_i$,
then
\[
 \llbracket\flatof{\Q(C;t_1,\ldots,t_k)}\rrbracket^{\N}_s
   =\code{C[\bn{n_1}/\h_1,\ldots,\bn{n_k}/\h_k]}.
\]
With $\atsign=s$, the prefix evaluation convention of
Section~\ref{at-sign-pairing-calculus} writes the same identity as
\[
 \atsign\,\flatof{\Q(C;t_1,\ldots,t_k)}
 =\code{C[\overline{\atsign\,t_1^{\flat}}/\h_1,\ldots,
                \overline{\atsign\,t_k^{\flat}}/\h_k]}.
\]
Here each $\atsign\,t_i^{\flat}$ is a number; its bar constructs a
numeral, which fills a hole before the resulting expression is coded.
These are semantic identities; \eqref{eq:lemma} is the separate
provability assertion for each fixed tuple of numerals. The use of
$\PA$ specifies the language in which the displayed function terms
exist. Their occurrences can be eliminated to obtain statements in PA.

\subsection{Occurrence profiles and arithmetized substitution}
\label{app:occurrence-feferman}

The finite occurrence construction of
Section~\ref{occurrence-indexed-substitution} applies to the term holes of
an arithmetic template. Each hole receives a term. The occurrence map
records which argument fills each slot, while the surrounding syntax
specifies the constructors of the resulting expression.

Let $C$ be a term or formula template with distinct holes
$\h_1,\ldots,\h_k$ and $m$ hole occurrences. Number the occurrences from
left to right and define
\[
 o_C:\finord m\longrightarrow\finord k,
 \qquad o_C(r)=i
 \quad\Longleftrightarrow\quad
 \text{the $r$-th hole occurrence is $\h_i$}.
\]
Let $C_{\mathrm{lin}}$ be the template obtained by replacing the $r$-th
hole occurrence by $\h_r$, for every $r\in\finord m$. Thus each hole of
$C_{\mathrm{lin}}$ occurs exactly once. This replacement is performed on
the occurrence slots of $C$ in one step. It leaves every ordinary variable,
function symbol, connective, and quantifier unchanged.

For any set $B$, precomposition with $o_C$ defines
\[
 o_C^*:B^{\finord k}\longrightarrow B^{\finord m},
 \qquad
 o_C^*(b_1,\ldots,b_k)
   =(b_{o_C(1)},\ldots,b_{o_C(m)}).
\]
Let $K\subseteq\N$ be the set of codes of closed $L^{*}$-terms. Choose
the fresh variables $w_1,\ldots,w_k$ used in the translation of $C$ and
write $C^{\mathbf w}$ as in \eqref{eq:grammar}. Define
\[
 b_C:K^{\finord k}\longrightarrow\N,
 \qquad
 b_C(\vec a)=
 \sbop_k\bigl(\code{C^{\mathbf w}};
       \code{w_1},\ldots,\code{w_k};a_1,\ldots,a_k\bigr)
 \quad(k>0).
\]
For $k=0$, the domain is a singleton and $b_C()=\code C$.
Define $b_{C_{\mathrm{lin}}}$ in the same way, using fresh variables
$z_1,\ldots,z_m$ for its holes. The variables in each list occur nowhere
in the corresponding template.

\begin{proposition}[Occurrence factorization]\label{prop:occurrence-feferman}
For every $\vec a\in K^{\finord k}$,
\begin{equation}\label{eq:occurrence-feferman}
 b_C(\vec a)=b_{C_{\mathrm{lin}}}(o_C^*\vec a).
\end{equation}
Consequently, the following diagram commutes, where the upper vertical
maps apply $\num$ coordinatewise:
\[
\begin{tikzcd}[column sep=huge,row sep=large]
 \N^{\finord k} \arrow[r,"o_C^*"]
   \arrow[d,"\num^{\finord k}"'] &
 \N^{\finord m} \arrow[d,"\num^{\finord m}"]\\
 K^{\finord k} \arrow[r,"o_C^*"] \arrow[d,"b_C"'] &
 K^{\finord m} \arrow[d,"b_{C_{\mathrm{lin}}}"]\\
 \N \arrow[r,equal] & \N.
\end{tikzcd}
\]
\end{proposition}

\begin{proof}
Write $a_i=\code{u_i}$ with $u_i$ a closed term. Substitution into
$C^{\mathbf w}$ inserts $u_{o_C(r)}$ at the $r$-th original hole
occurrence. Substitution into $C_{\mathrm{lin}}^{\mathbf z}$ with the
argument list $o_C^*(u_1,\ldots,u_k)$ inserts the same term at the same
position. Structural recursion on $C$ therefore yields the raw-syntax
identity
\[
 C^{\mathbf w}[u_1/w_1,\ldots,u_k/w_k]
 =C_{\mathrm{lin}}^{\mathbf z}
     [u_{o_C(1)}/z_1,\ldots,u_{o_C(m)}/z_m].
\]
All inserted terms are closed, so no variable can be captured and the
specified substitution procedure preserves the bound-variable names.
Coding this identity proves \eqref{eq:occurrence-feferman}. The upper
square commutes coordinatewise: both paths send $\vec n$ to
$(\num(n_{o_C(r)}))_{r\in\finord m}$. When $k=0$, also $m=0$,
and both lower maps have the constant value $\code C$.
\end{proof}

For surface arithmetic terms $t_1,\ldots,t_k$ and a current assignment
$\atsign$ in the standard arithmetic structure, the diagram and the
translation \eqref{eq:grammar} yield
\begin{equation}\label{eq:occurrence-feferman-at}
 \atsign\,\flatof{\Q(C;t_1,\ldots,t_k)}
 =\atsign\,\flatof{
       \Q(C_{\mathrm{lin}};t_{o_C(1)},\ldots,t_{o_C(m)})}.
\end{equation}
For $m>0$, the arithmetic term evaluated on the right expands to
\[
 \Sb_m\bigl(\quo{C_{\mathrm{lin}}^{\mathbf z}};
   \quo{z_1},\ldots,\quo{z_m};
   \Num(t_{o_C(1)}^{\flat}),\ldots,
   \Num(t_{o_C(m)}^{\flat})\bigr).
\]
Thus the occurrence map duplicates and reorders the represented numeral
arguments of Feferman's substitution operation. For each fixed tuple of
standard numerals, the equality between the two closed translated terms
is provable in $\PA$ by their defining computations, as in
Proposition~\ref{prop:evaluation}.

\begin{example}[A repeated numeral argument]
For $C=\h_1+\h_1$, the profile is $o_C=(1,1)$ and
$C_{\mathrm{lin}}=\h_1+\h_2$. The upper horizontal map in the diagram
is the diagonal $n\mapsto(n,n)$. In the prefix notation,
\[
 \atsign\,\flatof{\Q(\h_1+\h_1;x)}
   =\code{\overline{\atsign\,x}+\overline{\atsign\,x}}.
\]
At $\atsign\,x=0$, the value is $\code{0+0}=33$ for the coding of
Appendix~\ref{app:numerals}, whereas $\num(0+0)=\code0=1$.
The code retains the addition symbol. The template $\h_1\cdot\h_1$
has the same occurrence profile and a different constructor, so its
substitution codes are different. The profile and the preserved syntax
together determine the operation.
\end{example}

\subsubsection*{Composition of templates}
For each $i\in\finord k$, let $D_i$ be a term template on the holes
$\mathsf g_{i,1},\ldots,\mathsf g_{i,\ell_i}$, with no ordinary
variables. Assume these hole families are pairwise disjoint. Let $D_i$
have $n_i$ hole occurrences and profile
$o_i:\finord{n_i}\to\finord{\ell_i}$.
Replace each $\h_i$ in $C$ by $D_i$ and denote the resulting template
by $C[\vec D]$. Index its holes in the block order of
Section~\ref{occurrence-indexed-substitution}.

Each hole occurrence in $C[\vec D]$ has coordinates $(r,t)$, where $r$
selects a hole occurrence of $C$ and $t$ selects an occurrence in the
inserted copy of $D_{o_C(r)}$. Its hole label is therefore recorded by
exactly the fibrewise map of that section:
\[
\begin{aligned}
 D_{C,\vec n}
   &=\bigoplus_{r=1}^{m}\finord{n_{o_C(r)}},
 &C_{\vec\ell}&=\bigoplus_{i=1}^{k}\finord{\ell_i},\\
 \gamma_{C,\vec o}:D_{C,\vec n}&\longrightarrow C_{\vec\ell},
 &\gamma_{C,\vec o}(r,t)&=
       \bigl(o_C(r),o_{o_C(r)}(t)\bigr).
\end{aligned}
\]
An assignment of closed-term codes $a:C_{\vec\ell}\to K$ pulls back
along $\gamma_{C,\vec o}$ to the occurrence labels
\[
 (a\circ\gamma_{C,\vec o})(r,t)
   =a\bigl(o_C(r),o_{o_C(r)}(t)\bigr).
\]
The ordinal encodings of Section~\ref{occurrence-indexed-substitution}
transport this map to $\delta_{C,\vec o}$.

Let $a_i=(a(i,1),\ldots,a(i,\ell_i))$. Filling $D_i$ with closed
terms produces a closed term, so $b_{D_i}(a_i)\in K$. The substitution
maps on codes form the typed commutative square
\[
\begin{tikzcd}[column sep=12em,row sep=large]
 K^{C_{\vec\ell}}
   \arrow[r,"{a\mapsto(b_{D_i}(a_i))_{i\in\finord k}}"]
   \arrow[d,"b_{C[\vec D]}"'] &
 K^{\finord k} \arrow[d,"b_C"]\\
 \N \arrow[r,equal] & \N.
\end{tikzcd}
\]
Indeed, both paths code the expression obtained by filling every final
hole of $C[\vec D]$ with the prescribed closed term. Filling the inner
templates first and then inserting them in $C$ yields that same raw
expression. No bound-variable renaming is required. Equivalently,
\begin{equation}\label{eq:feferman-template-composition}
 b_{C[\vec D]}(a)
   =b_C\bigl(b_{D_1}(a_1),\ldots,b_{D_k}(a_k)\bigr).
\end{equation}
The map $\gamma_{C,\vec o}$ specifies the occurrence routing for this
composition; the primitive recursive substitution functions compute its
output code. Numeral insertion into the final holes is the specialization
$a(i,s)=\num(n_{i,s})$.


\subsection{Three provability assertions}

Fix a recursively enumerable theory $T$ in $L^{*}$ whose nonlogical
axioms are sentences, and fix an arithmetized proof predicate that
allows open formulas as proof conclusions. Let $\Prov_T(a)$ assert
that some finite $T$-proof has conclusion with code $a$. Its standard
interpretation is the actual theoremhood relation of $T$.
Fix $\varphi(v)$ with $v$ as its only free variable.

\begin{center}
\renewcommand{\arraystretch}{1.35}
\begin{tabularx}{\linewidth}{@{}p{0.43\linewidth}X@{}}
\toprule
Closed arithmetic sentence & What its truth in $\N$ means\\
\midrule
$\Prov_T(\quo{\forall v\,\varphi(v)})$
  & $T\vdash\forall v\,\varphi(v)$.\\
$\Prov_T(\quo{\varphi(v)})$
  & $T\vdash\varphi(v)$, with $v$ free in the conclusion.\\
$\forall x\,\Prov_T(q_\varphi(x))$
  & For every $n\in\N$, $T\vdash\varphi(\bn n)$.\\
\bottomrule
\end{tabularx}
\end{center}

The first two express equivalent theoremhood conditions: universal
elimination proves one direction and generalization proves the other,
because $v$ is not free in any axiom. This hypothesis is explicit in
\cite[\S2(c), Proposition~2.1, p.~45]{F60}; see also
\cite[\S2(b), p.~266]{F62}.

The third concerns proofs of numeral instances. Precisely,
\begin{equation}\label{eq:pointwise}
 \N\models\forall x\,\Prov_T(q_\varphi(x))
 \quad\Longleftrightarrow\quad
 (\forall n\in\N)\;T\vdash\varphi(\bn n).
\end{equation}
This follows from the value of $q_\varphi(x)$ at $x=n$ and the
correctness of the standard proof predicate. Proposition~\ref{prop:evaluation} also
yields, for every fixed $n$,
\[
 \PA\vdash
 \Prov_T(q_\varphi(\bn n))
 \leftrightarrow \Prov_T(\quo{\varphi(\bn n)}).
\]
This last inference uses equality congruence in arithmetic. It is not
a reflection principle.

Equation~\eqref{eq:pointwise} does not say that $T$ proves the
displayed universal provability sentence, nor that $T$ proves
$\forall v\,\varphi(v)$. In particular, a family of proofs of
numeral instances is not a proof with a free variable to which the
generalization rule can be applied.

For example, if $E_{e,f}(v)$ is a fixed formula expressing the equality
under consideration, with $e,f$ fixed, the two relations can be defined by
\begin{align*}
 e\equiv_T f&\quad\Longleftrightarrow\quad
       T\vdash\forall v\,E_{e,f}(v),\\
 e\approx_T f&\quad\Longleftrightarrow\quad
       (\forall n\in\N)\;T\vdash E_{e,f}(\bn n).
\end{align*}
The first row of the table expresses the former and the third row the
latter, when interpreted in $\N$ with $\varphi=E_{e,f}$. The formula
chosen for equality, especially for partial functions, remains part of
these definitions.

\subsection{Nesting: coding a term and coding its value}

Quotation is applied to a fixed arithmetic expression after its inner
abbreviations have been expanded. Consequently
\begin{equation}\label{eq:nesting-closed}
 \quo{\Prov_T\bigl(\quo{\varphi(\dot x)}\bigr)}
   =\quo{\Prov_T(q_\varphi(x))}
\end{equation}
is a closed numeral: it codes an open formula. The free $x$ in the
inner arithmetic term is part of the syntax coded by the outer quotation.

Two different parameterized constructions are useful. Put
\[
 F(v)\defeq\Prov_T(q_\varphi(v)),\qquad
 G(v)\defeq\Prov_T(v),\qquad m_n\defeq\code{\varphi(\bn n)}.
\]
The $q$ notation always designates substitution for the indicated free
variable of the fixed, expanded formula. In particular,
$q_\varphi(v)$ in the definition of $F$ is its explicit arithmetic term.
Then
\begin{align}
 \bigl(q_F(x)\bigr)^{\N,x\mapsto n}
     &=\code{\Prov_T(q_\varphi(\bn n))},
       \label{eq:nest-one}\\
 \bigl(q_G(q_\varphi(x))\bigr)^{\N,x\mapsto n}
     &=\code{\Prov_T(\bn{m_n})}.
       \label{eq:nest-two}
\end{align}
The first construction inserts
$\bn n$ into the syntax of the computation $q_\varphi(v)$; the
second inserts the numeral for that computation's \emph{value} into $G$.
Explicitly, the term occurring inside the first coded formula is
\[
 q_\varphi(\bn n)
   =\Sb\bigl(\quo{\varphi(v)};\quo v;\Num(\bn n)\bigr),
\]
for the designated variable $v$ of $\varphi$. It is not the numeral
$\bn{m_n}$. These terms have different syntax. With the fixed proof
predicate, whose argument occurs in its expansion, the two coded
formulas also have different syntax and hence different codes.

On the other hand, Proposition~\ref{prop:evaluation} yields
\[
 \PA\vdash q_\varphi(\bn n)=\bn{m_n},
 \qquad
 \PA\vdash
 \Prov_T(q_\varphi(\bn n))\leftrightarrow\Prov_T(\bn{m_n}).
\]
Provable equality of terms licenses replacement in formulas; it does
not identify the codes of the terms or of the resulting formulas.
The same distinction already occurs between $\Num(\bn n)$ and
$\bn{\num(n)}$ in Appendix~\ref{app:numerals}.

For reading nested expressions, expand the numeral-constructor and
substitution abbreviations before taking the outer code. The terms
$q_F(x)$ and $q_G(q_\varphi(x))$ display the two constructions explicitly.
A notation assigning quotation depths to repeated dots would require
additional syntactic conventions.

\subsection{Axiom parameters and consistency formulas}

\subsubsection*{A formula and a function producing its code}
An \emph{axiom formula} $\alpha(u)$, with
$\FV(\alpha)\subseteq\{u\}$, defines a set of axiom codes by
\[
 A_\alpha\defeq\{a\in\N:\N\models\alpha(\bn a)\}.
\]
Assume it selects only codes of sentences. Fix a proof calculus and
an arithmetical formula $\Prf_\alpha(a,p)$ saying that $p$ codes a
finite derivation with conclusion code $a$, whose nonlogical axioms
satisfy $\alpha$. Choose this formula by a fixed primitive recursive
construction from the syntax of $\alpha$. Define
\begin{equation}\label{eq:consistency}
 \Prov_\alpha(a)\defeq\exists p\,\Prf_\alpha(a,p),\qquad
 \Con_\alpha\defeq\neg\Prov_\alpha(\quo\bot),
\end{equation}
where $\bot$ is the fixed logical contradiction
$(0=0)\land\neg(0=0)$. This choice also works for fragments not
containing arithmetic axioms. Feferman's formulation excludes a
formula and its negation both being provable; it is equivalent to
\eqref{eq:consistency} for this classical proof calculus
\cite[\S2(c), p.~269]{F62}.

The formula $\Con_\alpha$ is distinct from the primitive recursive
function taking $\code\alpha$ to $\code{\Con_\alpha}$. Relative to
the fixed construction in \eqref{eq:consistency}, denote that numerical
function by $\mathrm{con}^{*}$ and its represented function symbol by
$\Con^{*}$. Thus
\begin{equation}\label{eq:con-star}
 \mathrm{con}^{*}(\code\alpha)=\code{\Con_\alpha},
 \qquad
 \PA\vdash\Con^{*}(\quo\alpha)=\quo{\Con_\alpha}.
\end{equation}
Set the numerical function to $0$ on inputs outside its prescribed
syntactic domain. Feferman makes exactly the distinction between
consistency formulas and represented code-producing functions on
p.~269 of \cite{F62}. The output G\"odel numbers depend on the selected
definition of the consistency formula.

\subsubsection*{Adding a numeral instance as an axiom}
Fix a formula $\varphi(v)$ with no free variables other than $v$,
and define
\begin{equation}\label{eq:beta}
 \beta(z;u)\defeq\alpha(u)\lor u=q_\varphi(z).
\end{equation}
Choose $z$ distinct from $u$. Here $z$ is a parameter and $u$ is
the variable for axiom codes.
For each $n\in\N$, $\beta(\bn n;u)$ defines the axiom set
\[
  A_\alpha\cup\{\code{\varphi(\bn n)}\}.
\]
The parameterized proof and consistency formulas are formed while
leaving $z$ free; auxiliary proof variables are chosen fresh.
We write the latter as $\Con_{\beta(z;\cdot)}$.
This is a formula with parameter $z$, not a term. The construction is
the parameter convention of \cite[\S2(c), p.~270]{F62}.

One may abbreviate it by $\Con_{\alpha+\varphi(\dot z)}$, provided
that the subscript is defined by \eqref{eq:beta}. The plus sign then
denotes adjoining an axiom through its axiom formula. It does not
assert that $\alpha+\varphi(\dot z)$ is itself an arithmetic term.
The arithmetic term selecting the added axiom is $q_\varphi(z)$;
a represented constructor such as $\Con^{*}$ is needed to compute
the code of a consistency formula.

\subsubsection*{The Di Paola--Montagna instance}
In the Di Paola--Montagna setting, $T$ is a consistent, recursively
enumerable extension of PA. Their characterization on p.~648 further
assumes that $T$ is reflexive and $\alpha$ is an RE-formula binumerating
its axioms in Robinson arithmetic. Let $\alpha$ be this axiom
formula, and choose the variables $z,x,u$ pairwise distinct. Define the bounded axiom formula
\begin{equation}\label{eq:restriction}
 (\alpha\upharpoonright z)(u)\defeq\alpha(u)\land u\leq z.
\end{equation}
Fix an index $e$ for a unary partial recursive function and a primitive
recursive computation predicate $\mathcal T(e,v,y)$ for which
$\varphi_e(v)$ is defined exactly when some $y$ satisfies
$\mathcal T(e,v,y)$. Let
\begin{align}
 \theta_e(v)&\defeq\exists y\,\mathcal T(\bn e,v,y),
                   \label{eq:theta}\\
 q_{\theta_e}(x)&\defeq
  \Sb\bigl(\quo{\theta_e(v)};\quo v;\Num(x)\bigr),
                   \label{eq:qtheta}\\
 \beta_e(z,x;u)&\defeq
    \bigl(\alpha(u)\land u\leq z\bigr)\lor u=q_{\theta_e}(x).
                   \label{eq:beta-e}
\end{align}
Thus $\theta_e(v)$ is the formula abbreviated by
$\varphi_e(v)\downarrow$ \cite[p.~645]{DPM91}. The dot in the added-axiom notation
specifies substitution of the numeral for $x$ into this fixed formula.

With the parameter arguments ordered as in
\eqref{eq:beta-e}, Di Paola and Montagna's abbreviation is
\begin{equation}\label{eq:dpm}
 \Con_{\alpha\upharpoonright\bn n+\varphi_e(\dot x)\downarrow}
   \defeq\Con_{\beta_e(\bn n,x;\cdot)}.
\end{equation}
The restriction uses $u\leq z$, including the endpoint. The bounded
axiom numeration and the consistency abbreviation occur on p.~646 of
\cite{DPM91}. Their condition~(ii), p.~648, has the quantifier pattern
\begin{equation}\label{eq:dpm-ii}
 (\forall n\in\N)\qquad
 T\vdash\forall x\,
   \Con_{\alpha\upharpoonright\bn n+\varphi_e(\dot x)\downarrow}.
\end{equation}
The quantification over $n$ is external: a separate theorem is asserted
for each standard finite bound. The quantifier $\forall x$ is inside
each theorem. For each fixed $n$, the formula concerns consistency of
the bounded fragment together with one convergence instance, as that
instance varies with $x$.

{\footnotesize
\begin{thebibliography}{99}
\setlength{\itemsep}{0pt}
\setlength{\parsep}{0pt}

\bibitem{DPcoherence}
Kosta Do\v{s}en and Zoran Petri\'{c},
\emph{Proof-Theoretical Coherence}, King's College Publications,
London, 2004.
\href{https://www.mi.sanu.ac.rs/novi_sajt/biography/kosta/coh.pdf}
{Revised version, September 2007}.

\bibitem{MacLane}
Saunders Mac Lane,
\emph{Categories for the Working Mathematician}, second edition,
Graduate Texts in Mathematics 5, Springer, 1998.

\bibitem{Lawvere1963}
F.~William Lawvere,
``Functorial semantics of algebraic theories,''
\emph{Proceedings of the National Academy of Sciences of the United States of
America} \textbf{50} (1963), 869--872.
\href{https://doi.org/10.1073/pnas.50.5.869}{doi:10.1073/pnas.50.5.869}.

\bibitem{CaramelloClassifying}
Olivia Caramello,
\emph{Topos Theory: Lectures 21 and 22---Classifying Toposes}.
\href{https://www.oliviacaramello.com/Teaching/CambridgeToposTheoryCourseLectures21and22.pdf}{Lecture notes}.

\bibitem{Manuell2022}
Graham Manuell,
``The spectrum of a localic semiring,''
\emph{Mathematical Proceedings of the Cambridge Philosophical Society}
\textbf{173} (2022), 647--668.
\href{https://doi.org/10.1017/S0305004122000068}{doi:10.1017/S0305004122000068}.

\bibitem{StacksPoints}
The Stacks Project Authors,
``Points of schemes,'' \emph{The Stacks Project}.
\href{https://stacks.math.columbia.edu/tag/01J5}{Tag 01J5}.

\bibitem{LurieUltraproducts}
Jacob Lurie,
\emph{Lecture 20X: Ultraproducts}, March 31, 2018.
\href{https://www.math.ias.edu/~lurie/278xnotes/Lecture20X-Ultraproducts.pdf}
{Lecture notes}.

\bibitem{HaineCondensedPoints}
Peter J. Haine,
\emph{Classifying anima of condensed \(\infty\)-categories of points},
\href{https://arxiv.org/abs/2602.21330}{arXiv:2602.21330}, 2026.

\bibitem{ScholzeCondensed}
Peter Scholze,
\emph{Lectures on Condensed Mathematics},
\href{https://arxiv.org/abs/2605.03658}{arXiv:2605.03658}, 2026.
Updated version of the 2019 lectures, presenting joint work with
Dustin Clausen.

\bibitem{F60}
Solomon Feferman.
\emph{Arithmetization of metamathematics in a general setting}.
Fundamenta Mathematicae \textbf{49} (1960), 35--92.
\href{https://doi.org/10.4064/fm-49-1-35-92}{doi:10.4064/fm-49-1-35-92}.

\bibitem{F62}
Solomon Feferman.
\emph{Transfinite recursive progressions of axiomatic theories}.
The Journal of Symbolic Logic \textbf{27}, no.~3 (1962), 259--316.
\href{https://doi.org/10.2307/2964649}{doi:10.2307/2964649}.

\bibitem{DPM91}
Robert A. Di Paola and Franco Montagna.
\emph{Some properties of the syntactic $p$-recursion categories generated
by consistent, recursively enumerable extensions of Peano arithmetic}.
The Journal of Symbolic Logic \textbf{56}, no.~2 (1991), 643--660.
\href{https://doi.org/10.2307/2274707}{doi:10.2307/2274707}.
\end{thebibliography}
}

\end{document}
